% 文件开头添加条件判断
\newif\ifstandalone
\standalonetrue  % 默认独立编译模式

% 检测是否在主文档中
\ifdefined\maindoc
  \standalonefalse
\fi

\ifstandalone
  \documentclass[UTF8]{ctexart}

  \usepackage{amsmath,amssymb,amsfonts}
  \usepackage{geometry}
  \usepackage{graphicx}
  \usepackage{booktabs}
  \usepackage{array}
  \usepackage{multirow}
  \usepackage{float}
  \usepackage{bm}

  \geometry{a4paper,margin=2.5cm}
\fi

\ifstandalone
  \begin{document}
  \section*{第七章\quad 先进控制技术}
\fi
% 主文档中章标题由 \chapter{先进控制技术} 生成，此处不再输出重复标题（否则目录无小节）

本章讨论机器人机械臂跟踪控制的一些先进控制技术。本章开发的控制器涉及计算问题和执行器动力学效应。本章提出的分析概念和控制开发通常比前几章提出的更为复杂；因此，强烈建议在研究新材料之前先学习前面的章节。

\section{引言}

随着过去几年来机器人控制研究的进展，许多机器人控制研究人员开始关注实现问题。也就是说，实现方面的考虑，如减少在线计算和执行器动力学效应，正在促使研究人员重新思考以前机器人控制器的理论开发，以便解决这些考虑。不断重新设计以前的控制开发以符合实现限制，是以前机器人控制研究进展的进行方式。利用这种迫使理论开发满足实现限制的概念，我们说明了研究人员如何开始着手解决诸如减少在线计算和补偿执行器动力学效应等问题。

\section{减少在线计算的机器人控制器}

在本节中，我们研究了由Sadegh及其合作者[Sadegh和Horowitz 1990]、[Sadegh等1990]设计的机器人控制器。我们将这些控制器与相关工作分开，因为这项工作解决了极其相关的在线控制器计算实现问题。具体而言，这种自适应控制器减少了在线计算，而不是其他控制技术，如第6章中提出的自适应控制器。在自适应控制器研究的发展之后，还提出了一种"重复"控制器。这种重复控制器也减少了在线计算。

\subsection{期望补偿自适应律}

第6章中自适应控制器的一个缺点是，用作前馈补偿的回归矩阵（例如，自适应惯量相关控制器中的矩阵$Y(\cdot)$）必须在线计算。回归矩阵必须在线计算，因为它取决于关节位置和速度的测量值（即$q$和$\dot{q}$）。对于例6.3.1中给出的简单二连杆机器人控制器，很明显$Y(\cdot)$的在线计算在计算上是密集的。可以想象，如果人们希望控制具有许多自由度的机器人机械臂，回归矩阵的在线计算可能在计算上非常密集。

为了消除对回归矩阵在线计算的需要，我们现在研究期望补偿自适应律（DCAL）[Sadegh和Horowitz 1990]。DCAL通过用期望关节位置和速度（即$q_d$和$\dot{q}_d$）替换$q$和$\dot{q}$，消除了对回归矩阵在线计算的需要。也就是说，DCAL回归矩阵仅取决于期望轨迹信息；因此，DCAL回归矩阵可以预先离线计算。当然，这种回归矩阵的修改迫使我们重新检查自适应控制设计和相应的稳定性分析。

为了本节中的控制设计，我们假设机器人机械臂是具有以下动力学的旋转机械臂：
\begin{equation}\label{eq:7.2.1}
\tau = M(q)\ddot{q} + V_m(q,\dot{q})\dot{q} + G(q) + F_d\dot{q}
\end{equation}
其中$F_d$是$n \times n$正定对角矩阵，用于表示摩擦的动态系数，所有其他量如第3章中定义。如其他章节所述，我们将关节跟踪误差定义为
\begin{equation}\label{eq:7.2.2}
e = q_d - q
\end{equation}

如第6章所述，机器人机械臂的自适应控制涉及将已知时间函数与未知常数参数分离。例如，回想自适应惯量相关控制器的这种参数与时间函数分离为
\begin{equation}\label{eq:7.2.3}
Y(\cdot)\varphi = M(q)(\ddot{q}_d + \dot{e}) + V_m(q,\dot{q})(\dot{q}_d + e) + G(q) + F_d\dot{q}
\end{equation}
其中$Y(\cdot)$是$n \times r$回归矩阵，仅取决于实际轨迹与期望轨迹的已知时间函数，$\varphi$是$r \times 1$未知常数参数向量。（注意表6.3.1中定义的$\Lambda$取为单位矩阵。）

在DCAL中，参数与时间函数的这种分离为
\begin{equation}\label{eq:7.2.4}
Y(\cdot)\varphi = M(q_d)\ddot{q}_d + V_m(q_d,\dot{q}_d)\dot{q}_d + G(q_d) + F_d\dot{q}_d
\end{equation}
其中式\eqref{eq:7.2.4}右端仅含期望轨迹$(q_d,\dot{q}_d,\ddot{q}_d)$；与式\eqref{eq:7.2.3}相同，左端记为$Y(\cdot)\varphi$，刻画同一未知参数向量$\varphi$。随后的DCAL扭矩及更新律中，将该期望轨迹回归矩阵记作$Y_d(\cdot)$以强调其与式\eqref{eq:7.2.3}中依赖于$(q,\dot{q})$等的$Y(\cdot)$的差别。注意，当$e=0$且$\dot{e}=0$（等价于$q=q_d$、$\dot{q}=\dot{q}_d$且$\ddot{q}=\ddot{q}_d$）时，式\eqref{eq:7.2.3}的右端化为式\eqref{eq:7.2.4}的右端。

利用式\eqref{eq:7.2.4}给出的回归矩阵公式，DCAL表示为
\begin{equation}\label{eq:7.2.5}
\tau = Y_d(\cdot)\hat{\varphi} + k_v r + k_p e + k_a \|e\|^2 r
\end{equation}
其中$k_v$、$k_p$、$k_a$是标量常数控制增益，$\hat{\varphi}$是$r \times 1$参数估计向量，滤波跟踪误差定义为
\begin{equation}\label{eq:7.2.6}
r = \dot{e} + e
\end{equation}

相应的DCAL参数自适应更新律为
\begin{equation}\label{eq:7.2.7}
\dot{\hat{\varphi}} = \Gamma Y_d^T(\cdot) r
\end{equation}
其中$\Gamma$是$r \times r$正定、对角、常数自适应增益矩阵，参数误差定义为
\begin{equation}\label{eq:7.2.8}
\tilde{\varphi} = \varphi - \hat{\varphi}
\end{equation}

注意，式\eqref{eq:7.2.5}给出的DCAL与第6章讨论的自适应控制器非常相似，除了式\eqref{eq:7.2.5}中的项$k_a\|e\|^2r$。事实证明，这个附加项用于补偿式\eqref{eq:7.2.3}和\eqref{eq:7.2.4}中给出的$Y(\cdot)\varphi$和$Y_d(\cdot)\varphi$之间的差异。如[Sadegh和Horowitz 1990]所示，实际回归矩阵与期望回归矩阵公式之间的这种差异可以量化为
\begin{equation}\label{eq:7.2.9}
\|\tilde{Y}\| \leq \zeta_1\|e\| + \zeta_2\|e\|^2 + \zeta_3\|r\| + \zeta_4\|r\|\|e\|
\end{equation}
其中
\begin{equation}\label{eq:7.2.10}
\tilde{Y} = Y(\cdot)\varphi - Y_d(\cdot)\varphi
\end{equation}
而$\zeta_1$、$\zeta_2$、$\zeta_3$和$\zeta_4$是正的边界常数，取决于期望轨迹和特定机器人配置的物理特性（即连杆质量、连杆长度、摩擦系数等）。

为了分析式\eqref{eq:7.2.5}给出的控制器的稳定性，我们必须形成相应的误差系统。首先，我们将式\eqref{eq:7.2.1}用式\eqref{eq:7.2.3}和\eqref{eq:7.2.6}中定义的$Y(\cdot)\varphi$和$r$重写。也就是说，我们有
\begin{equation}\label{eq:7.2.11}
M(q)\dot{r} = -V_m(q,\dot{q})r + Y(\cdot)\varphi - \tau
\end{equation}

在式\eqref{eq:7.2.11}的右侧添加和减去项$Y_d(\cdot)\varphi$得到
\begin{equation}\label{eq:7.2.12}
M(q)\dot{r} = -V_m(q,\dot{q})r + Y_d(\cdot)\varphi + \tilde{Y} - \tau
\end{equation}
其中$\tilde{Y}$在式\eqref{eq:7.2.10}中定义。将式\eqref{eq:7.2.5}给出的控制代入式\eqref{eq:7.2.12}得到误差系统
\begin{equation}\label{eq:7.2.13}
M(q)\dot{r} = -V_m(q,\dot{q})r - k_v r - k_p e - k_a\|e\|^2r + Y_d(\cdot)\tilde{\varphi} + \tilde{Y}
\end{equation}
其中$\tilde{\varphi}$在式\eqref{eq:7.2.8}中定义。

我们现在用Lyapunov-like函数分析式\eqref{eq:7.2.13}给出的误差系统的稳定性
\begin{equation}\label{eq:7.2.14}
V = \frac{1}{2}r^T M(q) r + \frac{1}{2}k_p e^T e + \frac{1}{2}\tilde{\varphi}^T \Gamma^{-1} \tilde{\varphi}
\end{equation}

对式\eqref{eq:7.2.14}关于时间求导得到
\begin{equation}\label{eq:7.2.15}
\dot{V} = \frac{1}{2}r^T\dot{M}(q)r + r^T M(q)\dot{r} + k_p e^T\dot{e} + \tilde{\varphi}^T \Gamma^{-1}\dot{\tilde{\varphi}}
\end{equation}
因为标量量可以转置。将式\eqref{eq:7.2.13}代入式\eqref{eq:7.2.15}得到
\begin{equation}\label{eq:7.2.16}
\begin{aligned}
\dot{V} = {} & k_p e^T\dot{e} - k_v r^T r - k_p r^T e - k_a\|e\|^2 r^T r + r^T\tilde{Y} \\
& \quad {} + \frac{1}{2} r^T \bigl( \dot{M}(q) - 2V_m(q,\dot{q}) \bigr) r + \tilde{\varphi}^T \Gamma^{-1}\dot{\tilde{\varphi}} + r^T Y_d(\cdot)\tilde{\varphi}
\end{aligned}
\end{equation}

通过利用斜对称性质（见第3章），式\eqref{eq:7.2.16}中$\frac{1}{2} r^T \bigl( \dot{M}(q) - 2V_m(q,\dot{q}) \bigr) r$项为零；再结合式\eqref{eq:7.2.7}的更新律可知$\tilde{\varphi}^T \Gamma^{-1}\dot{\tilde{\varphi}}+r^T Y_d(\cdot)\tilde{\varphi}=0$。因此，由式\eqref{eq:7.2.6}中$r$的定义，式\eqref{eq:7.2.16}简化为
\begin{equation}\label{eq:7.2.17}
\dot{V} = -k_p e^T e - k_v r^T r - k_a\|e\|^2 r^T r + r^T\tilde{Y}
\end{equation}

从式\eqref{eq:7.2.17}，我们可以用以下方式给出$\dot{V}$的上界：
\begin{equation}\label{eq:7.2.18}
\dot{V} \leq -k_p \|e\|^2 - k_v \|r\|^2 - k_a \|e\|^2 \|r\|^2 + \|r\|\|\tilde{Y}\|
\end{equation}

在式\eqref{eq:7.2.18}中代入式\eqref{eq:7.2.9}，可得$\dot{V}$的新上界
\begin{equation}\label{eq:7.2.19}
\dot{V} \leq -k_p \|e\|^2 - k_v \|r\|^2 - k_a \|e\|^2 \|r\|^2 + \zeta_1 \|e\| \|r\| + \zeta_2 \|e\|^2 \|r\| + \zeta_3 \|r\|^2 + \zeta_4 \|r\|^2 \|e\|
\end{equation}

将式\eqref{eq:7.2.19}右端配方后可写成
\begin{equation}\label{eq:7.2.20}
\begin{aligned}
\dot{V} \leq {} & -k_p \|e\|^2 - k_v \|r\|^2 - k_a \|e\|^2 \|r\|^2 + \zeta_1 \|e\| \|r\| \\
& \quad {} - \zeta_2 \|e\|^2 \bigl(\tfrac{1}{2} - \|r\|\bigr)^2 - \zeta_4 \|r\|^2 \bigl(\tfrac{1}{2} - \|e\|\bigr)^2 \\
& \quad {} + (\zeta_2 + \zeta_4) \|e\|^2 \|r\|^2 + \frac{\zeta_2}{4} \|e\|^2 + \Bigl(\zeta_3 + \frac{\zeta_4}{4}\Bigr) \|r\|^2
\end{aligned}
\end{equation}

在式\eqref{eq:7.2.20}中合并同类项后，它可以重写为
\begin{equation}\label{eq:7.2.21}
\begin{aligned}
\dot{V} \leq {} & -(k_p - \zeta_2/4) \|e\|^2 - (k_v - \zeta_3 - \zeta_4/4) \|r\|^2 + \zeta_1 \|e\| \|r\| \\
& \quad {} - \zeta_2 \|e\|^2 \bigl(\tfrac{1}{2} - \|r\|\bigr)^2 - \zeta_4 \|r\|^2 \bigl(\tfrac{1}{2} - \|e\|\bigr)^2 \\
& \quad {} - (k_a - \zeta_2 - \zeta_4) \|e\|^2 \|r\|^2
\end{aligned}
\end{equation}

注意到，若控制增益$k_a$按下式选取，
\begin{equation}\label{eq:7.2.22}
k_a > \zeta_2 + \zeta_4
\end{equation}
则式\eqref{eq:7.2.21}第三行非正；又因含$\bigl(\tfrac{1}{2}-\|r\|\bigr)^2$与$\bigl(\tfrac{1}{2}-\|e\|\bigr)^2$的两项非正，可得$\dot{V}$的上界
\begin{equation}\label{eq:7.2.23}
\dot{V} \leq -(k_p - \zeta_2/4) \|e\|^2 - (k_v - \zeta_3 - \zeta_4/4) \|r\|^2 + \zeta_1 \|e\| \|r\|
\end{equation}

通过以矩阵形式重写式\eqref{eq:7.2.23}
\begin{equation}\label{eq:7.2.24}
\dot{V} \leq -x^T \bar{Q} x
\end{equation}
其中
\begin{align*}
x &= \begin{bmatrix} \|e\| \\ \|r\| \end{bmatrix}, \\
\bar{Q} &= \begin{bmatrix} k_p - \zeta_2/4 & -\zeta_1/2 \\ -\zeta_1/2 & k_v - \zeta_3 - \zeta_4/4 \end{bmatrix}.
\end{align*}

我们可以建立$k_p$和$k_v$的充分条件，使得$\bar{Q}$为正定。具体而言，通过使用Gerschgorin定理（见第2章），我们可以看到，如果
\begin{equation}\label{eq:7.2.25}
k_p > \zeta_1/2 + \zeta_2/4
\end{equation}
和
\begin{equation}\label{eq:7.2.26}
k_v > \zeta_1/2 + \zeta_3 + \zeta_4/4
\end{equation}
$\bar{Q}$将为正定；因此，$\dot{V}$将为半负定。

我们现在详细说明跟踪误差的稳定性类型。首先，由于$\dot{V}$是半负定的，我们可以说明$V$是上界有界的。利用$V$是上界有界的事实，我们可以说明$e$、$\dot{e}$、$r$和$\tilde{\varphi}$是有界的。由于$e$、$\dot{e}$、$r$和$\tilde{\varphi}$是有界的，我们可以使用式\eqref{eq:7.2.13}来表明$\dot{r}$以及因此式\eqref{eq:7.2.17}中的$\ddot{V}$是有界的。其次，注意由于$M(q)$如惯性矩阵的正定性质所阐述的那样下界有界（见第3章），我们可以说明式\eqref{eq:7.2.14}中给出的$V$下界有界。由于$V$下界有界，$\dot{V}$是半负定的，且$\ddot{V}$是有界的，我们可以使用Barbalat引理（见第2章）来说明
\begin{equation*}
\lim_{t \to \infty} \dot{V} = 0
\end{equation*}

因此，从上面的论证和式\eqref{eq:7.2.24}，我们知道
\begin{equation}\label{eq:7.2.27}
\lim_{t \to \infty} \begin{bmatrix} e \\ r \end{bmatrix} = 0
\end{equation}

从式\eqref{eq:7.2.27}，我们也可以确定速度跟踪误差的稳定性结果。具体而言，从式\eqref{eq:7.2.6}，注意$r$被定义为变量$e$的稳定一阶微分方程；因此，通过标准线性控制论证，我们可以写成
\begin{equation}\label{eq:7.2.28}
\lim_{t \to \infty} \|\dot{e}\| = 0
\end{equation}

这个结果告诉我们，如果控制器增益根据式\eqref{eq:7.2.22}、\eqref{eq:7.2.25}和\eqref{eq:7.2.26}选择，跟踪误差$e$和$\dot{e}$是渐近稳定的。从上面的分析，关于参数误差我们所能说的就是它保持有界。刚刚推导的自适应控制器总结在表7.2.1中，并在图7.2.1中描绘。

\textbf{表7.2.1}\quad DCAL控制器

\begin{center}
\begin{tabular}{|l|l|}
\hline
\textbf{控制律} & $\tau = Y_d(\cdot)\hat{\varphi} + k_v r + k_p e + k_a \|e\|^2 r$ \\
\hline
\textbf{滤波跟踪误差} & $r = \dot{e} + e$ \\
\hline
\textbf{参数更新律} & $\dot{\hat{\varphi}} = \Gamma Y_d^T(\cdot) r$ \\
\hline
\end{tabular}
\end{center}

浏览表7.2.1后，我们可以看到，与自适应惯量相关控制器相比，DCAL具有明显的减少在线计算的优势。具体而言，回归矩阵$Y_d(\cdot)$仅取决于期望轨迹；因此，回归矩阵可以离线计算。我们现在提出一个示例来说明如何使用表7.2.1为机器人机械臂设计自适应控制器。

\textbf{示例7.2--1：二连杆机械臂的DCAL}

我们希望为图6.2.1中给出的二连杆机械臂设计和仿真表7.2.1中给出的DCAL。（该机器人臂的动力学在第3章中给出。）假设摩擦可忽略不计且连杆长度精确知道为每根1米长，DCAL可以写成
\begin{align}
\tau_1 &= Y_{11}\hat{\varphi}_1 + Y_{12}\hat{\varphi}_2 + k_v r_1 + k_p e_1 + k_a \|e\|^2 r_1, \\
\tau_2 &= Y_{21}\hat{\varphi}_1 + Y_{22}\hat{\varphi}_2 + k_v r_2 + k_p e_2 + k_a \|e\|^2 r_2
\end{align}

在控制扭矩的表达式中，回归矩阵$Y_d(\ddot{q}_d,\dot{q}_d,q_d)$由
\begin{equation}
Y_d(\ddot{q}_d,\dot{q}_d,q_d) = \begin{bmatrix}
Y_{11} & Y_{12} \\
Y_{21} & Y_{22}
\end{bmatrix}
\end{equation}
给出，其中
\begin{align}
Y_{11} &= l_1^2 \ddot{q}_{d1} + l_1 g\cos(q_{d1}), \\
\begin{split}
Y_{12} &= (l_2^2 + 2l_1l_2\cos(q_{d2}) + l_1^2)\ddot{q}_{d1} + (l_2^2 + l_1l_2\cos(q_{d2}))\ddot{q}_{d2} \\
& \quad {} - l_1l_2\sin(q_{d2})\dot{q}_{d2}\dot{q}_{d1} - l_1l_2\sin(q_{d2})(\dot{q}_{d1}+\dot{q}_{d2})\dot{q}_{d2} \\
& \quad {} + l_2 g\cos(q_{d1}+q_{d2}) + l_1 g\cos(q_{d1}),
\end{split} \\
Y_{21} &= 0, \\
Y_{22} &= (l_1l_2\cos(q_{d2}) + l_2^2)\ddot{q}_{d1} + l_2^2 \ddot{q}_{d2} + l_1l_2\sin(q_{d2})\dot{q}_{d1}^2.
\end{align}

按照表7.2.1给出的公式制定自适应更新规则，相关参数估计向量为
\begin{equation*}
\hat{\varphi} = \begin{bmatrix} \hat{m}_1 \\ \hat{m}_2 \end{bmatrix}
\end{equation*}
带有自适应更新规则
\begin{align}
\dot{\hat{m}}_1 &= \gamma_1(Y_{11}r_1 + Y_{21}r_2) = \gamma_1 Y_{11} r_1, \\
\dot{\hat{m}}_2 &= \gamma_2(Y_{12}r_1 + Y_{22}r_2)
\end{align}

对于$m_1=0.8$ kg和$m_2=2.3$ kg，DCAL在以下条件下进行了仿真
\begin{equation*}
k_v = k_p = k_a = 50, \quad \gamma_1 = \gamma_2 = 20
\end{equation*}
和
\begin{equation*}
q_{d1} = q_{d2} = \sin t, \quad \dot{q}_{d1} = \dot{q}_{d2} = \cos t
\end{equation*}
\begin{equation*}
\hat{m}_1(0) = 0, \quad \hat{m}_2(0) = 0
\end{equation*}

跟踪误差和质量估计在图7.2.2中描绘。如图所示，跟踪误差是渐近稳定的，参数估计保持有界。

\subsection{重复控制律}

在许多工业应用中，机器人机械臂被用来重复执行相同的任务。例如，机器人可能被要求在装配线上反复喷涂相同的零件。可以想象，这种喷涂操作的期望轨迹将是周期函数。也就是说，在为喷涂第一个零件生成期望轨迹之后（即第一个"试验"），相同的期望轨迹应该以重复的方式为喷涂下一个零件而遵循（即下一个"试验"）。

如果控制策略不考虑重复操作的性质，沿第一条轨迹犯的错误将从试验到试验重复。因此，人们有动机设计一种控制器，利用当前试验中的跟踪误差测量来改善下一次试验的跟踪性能。这些类型的控制器通常被称为"学习"或重复控制器。"学习控制器"一词用于强调控制器试图学习机械臂动力学的可重复部分。

为了推动重复控制律（RCL）[Sadegh等1990]的设计，我们注意到由
\begin{equation}\label{eq:7.2.29}
u_d(t) = M(q_d)\ddot{q}_d + V_m(q_d,\dot{q}_d)\dot{q}_d + G(q_d) + F_d\dot{q}_d
\end{equation}
给出的动力学如果期望轨迹是周期性的，则是可重复的。也就是说，即使式\eqref{eq:7.2.29}中可能存在未知的常数参数，由$n \times 1$向量$u_d(t)$表示的信号将是周期性的或可重复的。因此，在随后的讨论中，我们假设期望轨迹的周期为$T$。期望轨迹的这种周期性假设使我们能够写成
\begin{equation}\label{eq:7.2.30}
u_d(t) = u_d(t-T)
\end{equation}
因为$u_d(t)$表示的动力学仅取决于周期性量。

利用式\eqref{eq:7.2.29}给出的动力学的可重复性，RCL表示为
\begin{equation}\label{eq:7.2.31}
\tau = \hat{u}_d(t) + k_v r + k_p e + k_a\|e\|^2r
\end{equation}
其中$n \times 1$向量$\hat{u}_d(t)$是用于补偿可重复动力学$u_d(t)$的学习项，所有其他量与DCAL中定义的相同。学习项$\hat{u}_d(t)$通过以下学习更新规则从试验到试验更新
\begin{equation}\label{eq:7.2.32}
\hat{u}_d(t) = \hat{u}_d(t-T) + k_L r
\end{equation}
其中$k_L$是正的标量控制增益。

如自适应控制开发中类似地进行，我们将用式\eqref{eq:7.2.32}给出的学习更新规则用学习误差表示，学习误差定义为
\begin{equation}\label{eq:7.2.33}
\tilde{u}_d(t) = u_d(t) - \hat{u}_d(t)
\end{equation}

具体而言，将式\eqref{eq:7.2.32}两边乘以$-1$，再在两边加上$u_d(t)$，可得
\begin{equation}\label{eq:7.2.34}
u_d(t) - \hat{u}_d(t) = u_d(t) - \hat{u}_d(t-T) - k_L r
\end{equation}

通过利用式\eqref{eq:7.2.30}给出的周期性假设，可将式\eqref{eq:7.2.34}写成
\begin{equation}\label{eq:7.2.35}
u_d(t) - \hat{u}_d(t) = u_d(t-T) - \hat{u}_d(t-T) - k_L r
\end{equation}
由此得到学习误差更新规则
\begin{equation}\label{eq:7.2.36}
\tilde{u}_d(t) = \tilde{u}_d(t-T) - k_L r
\end{equation}
其中$\tilde{u}_d(t-T)=u_d(t-T)-\hat{u}_d(t-T)$（对照式\eqref{eq:7.2.33}）。

在分析式\eqref{eq:7.2.31}中给出的控制器的稳定性之前，我们将形成相应的误差系统。首先，我们将式\eqref{eq:7.2.1}用式\eqref{eq:7.2.6}中定义的$r$重写。也就是说，我们有
\begin{equation}\label{eq:7.2.37}
M(q)\dot{r} + V_m(q,\dot{q})r + u_a(t) = \tau
\end{equation}
其中$n \times 1$向量$u_a(t)$用于表示"实际机械臂动力学"
\begin{equation}\label{eq:7.2.38}
u_a(t) = M(q)(\ddot{q}_d + \dot{e}) + V_m(q,\dot{q})(\dot{q}_d + e) + G(q) + F_d(\dot{q}_d + e)
\end{equation}

在式\eqref{eq:7.2.37}的右侧添加和减去项$u_d(t)$得到
\begin{equation}\label{eq:7.2.39}
M(q)\dot{r} + V_m(q,\dot{q})r + u_d(t) + \tilde{u} = \tau
\end{equation}
其中$\tilde{u}$定义为
\begin{equation}\label{eq:7.2.40}
\tilde{u} = u_a(t) - u_d(t)
\end{equation}

如[Sadegh和Horowitz 1990]中类似所示，实际机械臂动力学（即$u_a(t)$）与可重复机械臂动力学（即$u_d(t)$）之间的这种差异可以量化为
\begin{equation}\label{eq:7.2.41}
\|\tilde{u}\| \leq \zeta_1\|e\| + \zeta_2\|e\|\|r\| + \zeta_3\|r\| + \zeta_4
\end{equation}
其中$\zeta_1$、$\zeta_2$、$\zeta_3$和$\zeta_4$是正的边界常数，取决于期望轨迹和特定机器人配置的物理特性（即连杆质量、连杆长度、摩擦系数等）。

形成误差系统的最后一步是将式\eqref{eq:7.2.31}给出的控制代入式\eqref{eq:7.2.39}得到
\begin{equation}\label{eq:7.2.42}
M(q)\dot{r} = -V_m(q,\dot{q})r - k_v r - k_p e - k_a\|e\|^2r + \tilde{u}_d(t) - \tilde{u}
\end{equation}

我们现在用Lyapunov-like函数分析式\eqref{eq:7.2.42}给出的误差系统的稳定性
\begin{equation}\label{eq:7.2.43}
V = \frac{1}{2}r^T M(q)r + \frac{1}{2}k_p e^T e + \frac{1}{2k_L}\int_{t-T}^{t}\tilde{u}_d^T(\sigma)\tilde{u}_d(\sigma)d\sigma
\end{equation}

对式\eqref{eq:7.2.43}关于时间求导得到
\begin{equation}\label{eq:7.2.44}
\dot{V} = r^T M(q)\dot{r} + \frac{1}{2}r^T\dot{M}(q)r + k_p e^T\dot{e} + \frac{1}{2k_L}[\tilde{u}_d^T(t)\tilde{u}_d(t) - \tilde{u}_d^T(t-T)\tilde{u}_d(t-T)]
\end{equation}

将式\eqref{eq:7.2.42}给出的误差系统代入式\eqref{eq:7.2.44}得到
\begin{equation}\label{eq:7.2.45}
\begin{aligned}
\dot{V} = & -r^T V_m(q,\dot{q})r - k_v r^T r - k_p r^T e - k_a\|e\|^2r^T r \\
& + r^T\tilde{u}_d(t) - r^T\tilde{u} + \frac{1}{2}r^T\dot{M}(q)r + k_p e^T\dot{e} + \frac{1}{2k_L}[\tilde{u}_d^T(t)\tilde{u}_d(t) - \tilde{u}_d^T(t-T)\tilde{u}_d(t-T)]
\end{aligned}
\end{equation}

通过利用斜对称性质和学习误差更新律式\eqref{eq:7.2.36}，很容易表明式\eqref{eq:7.2.45}中的第二行等于
\begin{equation*}
-\frac{1}{2k_L}[\tilde{u}_d(t) + k_L r]^T[\tilde{u}_d(t) + k_L r]
\end{equation*}

因此，通过调用式\eqref{eq:7.2.6}中给出的$r$的定义，式\eqref{eq:7.2.45}简化为
\begin{equation}\label{eq:7.2.46}
\dot{V} = -k_v r^T r - k_p e^T e - k_a\|e\|^2r^T r + r^T\tilde{u}
\end{equation}

从式\eqref{eq:7.2.46}我们可以以以下方式放置$\dot{V}$的上界：
\begin{equation}\label{eq:7.2.47}
\dot{V} \leq -k_v \|r\|^2 - k_p \|e\|^2 - k_a\|e\|^2\|r\|^2 + \|r\|\|\tilde{u}\|
\end{equation}

其余稳定性论证是前一节对DCAL提出的稳定性论证的修改。具体而言，我们首先注意到式\eqref{eq:7.2.18}和\eqref{eq:7.2.47}几乎相同，因为式\eqref{eq:7.2.18}中的$\tilde{Y}$和式\eqref{eq:7.2.47}中的$\tilde{u}$由相同的标量函数界定。追溯DCAL稳定性论证的步骤后，我们可以看到控制器增益$k_a$和$k_p$仍应按照式\eqref{eq:7.2.22}和\eqref{eq:7.2.25}分别调整。然而，控制器增益$k_v$与控制器增益$k_L$一起调整以满足
\begin{equation}\label{eq:7.2.48}
k_v > \zeta_3 + \frac{1}{2} + k_L\zeta_2^2
\end{equation}
其中$\zeta_1$、$\zeta_2$、$\zeta_3$和$\zeta_4$在式\eqref{eq:7.2.41}中定义。如果控制器增益根据式\eqref{eq:7.2.22}、\eqref{eq:7.2.25}和\eqref{eq:7.2.48}调整，则从DCAL的分析开发（即式\eqref{eq:7.2.18}到\eqref{eq:7.2.24}）和式\eqref{eq:7.2.47}，我们可以在$\dot{V}$上放置新的上界：
\begin{equation}\label{eq:7.2.49}
\dot{V} \leq -\lambda_3 \|x\|^2
\end{equation}
其中$\lambda_3$为正的标量常数，且$\lambda_3=\lambda_{\min}\{Q_o\}$，而
\begin{align*}
x &= \begin{bmatrix} \|e\| \\ \|r\| \end{bmatrix}, \\
Q_o &= \begin{bmatrix} k_p - \zeta_2/4 & -\zeta_1/2 \\ -\zeta_1/2 & k_v + \dfrac{k_L}{2} - \zeta_3 - \zeta_4/4 \end{bmatrix}.
\end{align*}

我们现在详细说明跟踪误差的稳定性类型。首先注意，由式\eqref{eq:7.2.49}还可得到$\dot{V}$的更松上界：
\begin{equation}\label{eq:7.2.50}
\dot{V} \leq -\lambda_3 \|r\|^2
\end{equation}
这是因为$\|x\|^2=\|e\|^2+\|r\|^2\geq \|r\|^2$。式\eqref{eq:7.2.50}意味着
\begin{equation}\label{eq:7.2.51}
\int_0^{\infty} \dot{V}(\sigma)\,d\sigma \leq -\lambda_3 \int_0^{\infty} \|r(\sigma)\|^2\,d\sigma
\end{equation}

将式\eqref{eq:7.2.51}乘以$-1$，并对左侧关于$\sigma$积分，可得
\begin{equation}\label{eq:7.2.52}
V(0) - V(\infty) \geq \lambda_3 \int_0^{\infty} \|r(\sigma)\|^2\,d\sigma
\end{equation}

由于$\dot{V}$如式\eqref{eq:7.2.49}所述为半负定，故$V$为非增函数，因而由上界$V(0)$界定。又因惯性矩阵$M(q)$正定，式\eqref{eq:7.2.43}给出的$V$由零下界界定。由于$V$非增、上有界$V(0)$、下有界零，可将式\eqref{eq:7.2.52}写成
\begin{equation}\label{eq:7.2.53}
\lambda_3 \int_0^{\infty} \|r(\sigma)\|^2\,d\sigma < \infty
\end{equation}
或等价地
\begin{equation}\label{eq:7.2.54}
\sqrt{\int_0^{\infty} \|r(\sigma)\|^2\,d\sigma} < \infty
\end{equation}

式\eqref{eq:7.2.54}表明$r\in L_2^n$（见第2章），即滤波跟踪误差$r$在$L_2$意义下平方可积。

为了建立位置跟踪误差$e$的稳定性结果，我们建立位置跟踪误差与滤波跟踪误差$r$之间的传递函数关系。从式\eqref{eq:7.2.6}，我们有
\begin{equation}\label{eq:7.2.55}
e(s) = G(s)r(s)
\end{equation}
其中$s$是拉普拉斯变换变量，
\begin{equation}\label{eq:7.2.56}
G(s) = (sI + \Gamma)^{-1}
\end{equation}
而$I$是$n \times n$单位矩阵；$\Gamma$为与式\eqref{eq:7.2.6}中$r=\dot{e}+e$相对应的$n\times n$矩阵（此处$\Gamma=I$，故亦可写成$G(s)=(sI+I)^{-1}$）。由于$G(s)$是严格真、渐近稳定的传递函数且$r\in L_2^n$，我们可以使用第2章中的定理2.4.7来说明
\begin{equation}\label{eq:7.2.57}
\lim_{t \to \infty} e = 0
\end{equation}

因此，如果控制器增益根据式\eqref{eq:7.2.22}、\eqref{eq:7.2.25}和\eqref{eq:7.2.48}选择，位置跟踪误差$e$是渐近稳定的。根据本节提出的理论开发，关于速度跟踪误差$\dot{e}$我们只能说它是有界的。应该注意的是，如果式\eqref{eq:7.2.32}中的学习估计$\hat{u}_d(t)$"人为地"防止增长，我们可以得出速度跟踪误差是渐近稳定的结论[Sadegh等1990]。这种修改的稳定性证明是第6章中提出的自适应控制证明的直接应用。

本节研究的重复控制器总结在表7.2.2中，并在图7.2.3中描绘。

\textbf{表7.2.2}\quad RCL控制器

\begin{center}
\begin{tabular}{|l|l|}
\hline
\textbf{扭矩控制器} & $\tau = \hat{u}_d + k_v r + k_p e + k_a \|e\|^2 r$ \\
\hline
\textbf{滤波跟踪误差} & $r = e + \dot{e}$ \\
\hline
\textbf{学习更新律} & $\hat{u}_d(t) = \hat{u}_d(t-T) + k_L r$ \\
\hline
\textbf{稳定性} & 跟踪误差$e$渐近稳定；跟踪误差$\dot{e}$有界 \\
\hline
\textbf{注释} & 期望轨迹须为周期$T$的周期轨迹；增益$k_a$、$k_p$、$k_L$、$k_v$须选得足够大 \\
\hline
\end{tabular}
\end{center}

浏览表7.2.2后，我们可以看到RCL与需要形成回归型矩阵的自适应控制器相比，需要关于被控机器人的很少信息。RCL的另一个明显优势是它需要很少的在线计算。我们现在提出一个示例来说明如何使用表7.2.2为机器人机械臂设计重复控制器。

\textbf{示例7.2--2：二连杆机械臂的RCL}

我们希望为图6.2.1中给出的二连杆机械臂设计和仿真表7.2.2中给出的RCL。（该机器人臂的动力学在第3章中给出。）从表7.2.2，RCL可以写成
\begin{align}
\tau_1 &= \hat{u}_{d1} + k_v r_1 + k_p e_1 + k_a\|e\|^2r_1, \\
\tau_2 &= \hat{u}_{d2} + k_v r_2 + k_p e_2 + k_a\|e\|^2r_2
\end{align}

按照表7.2.2给出的公式制定学习更新规则得到
\begin{align}
\hat{u}_{d1}(t) &= \hat{u}_{d1}(t-T) + k_L r_1, \\
\hat{u}_{d2}(t) &= \hat{u}_{d2}(t-T) + k_L r_2
\end{align}

对于$m_1=0.8$ kg，$m_2=2.3$ kg，连杆长度为1 m，RCL在以下条件下进行了仿真
\begin{equation*}
k_a = k_v = k_p = k_L = 50, \quad T = 2\pi
\end{equation*}
和期望轨迹
\begin{equation*}
q_{d1} = \sin t, \quad q_{d2} = \sin t
\end{equation*}
以及初始条件
\begin{equation*}
q_1(0)=0,\quad q_2(0)=0,\quad \dot{q}_1(0)=0,\quad \dot{q}_2(0)=0,
\end{equation*}
学习项初值取
\begin{equation*}
\hat{u}_{d1}(t) = 0, \quad \hat{u}_{d2}(t) = 0, \quad \text{对于} -2\pi \leq t < 0
\end{equation*}

位置和速度跟踪误差在图7.2.4中描绘。如图所示，位置和速度跟踪误差都是渐近稳定的；然而，根据本小节的理论开发，我们只保证位置跟踪误差是渐近稳定的。

\section{自适应鲁棒控制}

在第6章中，我们讨论了自适应控制器用于机器人机械臂跟踪控制的用途。自适应控制器的有吸引力的特征之一是，控制实现不需要诸如负载质量或摩擦系数之类的未知常数参数的先验知识。自适应控制器的两个缺点是需要大量的在线计算，以及缺乏对附加有界扰动的鲁棒性。

在第5章中，我们讨论了鲁棒控制器用于机器人机械臂控制的用途。鲁棒控制器的两个有吸引力的特征是在线计算保持在最低限度，以及它们对附加有界扰动的固有鲁棒性。鲁棒控制方法的一个缺点是需要不确定性的先验已知界限。通常，不确定性界限的计算可能是一个相当繁琐的过程，因为此计算涉及为机器人机械臂的每个连杆找到质量和摩擦相关常数的最大值。鲁棒控制方法的另一个缺点是，即使在不存在附加有界扰动的情况下，我们也不能保证跟踪误差的渐近稳定性。通常，至少获得跟踪误差的"理论"渐近稳定性结果将是可取的。

在本节中，为机器人机械臂的跟踪控制开发了自适应鲁棒控制器。自适应鲁棒控制器可以被认为是结合了自适应控制器和鲁棒控制器的最佳品质。这种控制方法具有减少在线计算（与自适应控制方法相比）、对附加有界扰动的鲁棒性、无需系统不确定性的先验知识以及渐近跟踪误差性能的优点。

为了本节中的控制设计，我们假设机器人机械臂是具有以下动力学的旋转机械臂：
\begin{equation}\label{eq:7.3.1}
\tau = M(q)\ddot{q} + V_m(q,\dot{q})\dot{q} + G(q) + F_d\dot{q} + F_s(\dot{q}) + T_d
\end{equation}
其中$F_d$是用于表示摩擦动态系数的$n \times n$正定对角矩阵，$F_s(\dot{q})$是表示静态摩擦的$n \times 1$向量，$T_d$是表示未知有界扰动的$n \times 1$向量，所有其他量如第3章中定义。

自适应鲁棒控制器与第5章中讨论的鲁棒控制策略非常相似，因为使用辅助控制器来"界定"不确定性。回想第5章，鲁棒控制器通过使用由跟踪误差范数和正边界常数组成的标量函数来界定不确定性。例如，假设由
\begin{equation}\label{eq:7.3.2}
w = M(q)(\ddot{q}_d + \dot{e}) + V_m(q,\dot{q})(\dot{q}_d + e) + G(q) + F_d\dot{q} + F_s(\dot{q}) + T_d
\end{equation}
给出的动力学表示给定机器人控制器的不确定性。也就是说，式\eqref{eq:7.3.2}给出的动力学是不确定的，因为负载质量、摩擦系数和扰动不确切知道。然而，假设可使用正的标量函数$\rho$界定不确定性，如下所示：
\begin{equation}\label{eq:7.3.3}
\|w\| \leq \rho
\end{equation}

如[Dawson等1990]中所述，机器人机械臂的物理特性可用于表明式\eqref{eq:7.3.2}给出的动力学可界定为
\begin{equation}\label{eq:7.3.4}
\rho = \delta_0 + \delta_1\|\mathbf{e}\| + \delta_2\|\mathbf{e}\|^2
\end{equation}
其中
\begin{equation}\label{eq:7.3.5}
\mathbf{e} = \begin{bmatrix} e \\ \dot{e} \end{bmatrix}
\end{equation}
而$\delta_0$、$\delta_1$和$\delta_2$是正的边界常数，基于最大可能的负载质量、连杆质量、摩擦系数、扰动等（由[Dawson等1990]，总可适当选取这些常数使式\eqref{eq:7.3.4}给出的$\rho$满足式\eqref{eq:7.3.3}）。

通常，第5章中提出的鲁棒控制器要求式\eqref{eq:7.3.4}中定义的边界常数事先制定。本节将开发的自适应鲁棒控制器在机械臂运动时在线"学习"这些边界常数。也就是说，在控制实现中，我们不需要边界常数的知识；相反，我们只需要式\eqref{eq:7.3.4}中定义的边界常数的存在。

与[Corless和Leitmann 1983]中提出的一般开发类似，自适应鲁棒控制器具有以下形式
\begin{equation}\label{eq:7.3.6}
\tau = K_v r + v_R
\end{equation}
其中$K_v$是$n \times n$对角、正定矩阵，$r$（滤波跟踪误差）如式\eqref{eq:7.2.6}中定义，而$v_R$是表示辅助控制器的$n \times 1$向量。式\eqref{eq:7.3.6}中的辅助控制器$v_R$由
\begin{equation}\label{eq:7.3.7}
v_R = \frac{r\hat{\rho}^2}{\hat{\rho}\|r\| + \varepsilon}
\end{equation}
定义
其中
\begin{equation}\label{eq:7.3.8}
\dot{\varepsilon} = -k_\varepsilon\varepsilon
\end{equation}
$k_\varepsilon$是正的标量控制常数，$\hat{\rho}$定义为标量函数
\begin{equation}\label{eq:7.3.9}
\hat{\rho} = \hat{\delta}_0 + \hat{\delta}_1\|\mathbf{e}\| + \hat{\delta}_2\|\mathbf{e}\|^2
\end{equation}
其中$\hat{\delta}_0$、$\hat{\delta}_1$和$\hat{\delta}_2$分别是$\delta_0$、$\delta_1$和$\delta_2$的动态估计（见式\eqref{eq:7.3.4}）。带"$\hat{}$"的边界估计基于自适应更新律在线更改。在给出更新律之前，我们将式\eqref{eq:7.3.9}写成更为方便的矩阵形式
\begin{equation}\label{eq:7.3.10}
\hat{\rho} = S\hat{\theta}
\end{equation}
其中
\begin{equation*}
S = \begin{bmatrix} 1 & \|\mathbf{e}\| & \|\mathbf{e}\|^2 \end{bmatrix}, \qquad
\hat{\theta} = \begin{bmatrix} \hat{\delta}_0 \\ \hat{\delta}_1 \\ \hat{\delta}_2 \end{bmatrix}
\end{equation*}

式\eqref{eq:7.3.4}中给出的实际边界函数$\rho$也可以写成矩阵形式
\begin{equation}\label{eq:7.3.11}
\rho = S\theta
\end{equation}
其中
\begin{equation*}
\theta = \begin{bmatrix} \delta_0 \\ \delta_1 \\ \delta_2 \end{bmatrix}
\end{equation*}

注意自适应方法中回归矩阵公式（见第6章）与式\eqref{eq:7.3.10}给出的公式之间的相似性。具体而言，$1 \times 3$矩阵$S$类似于"回归矩阵"，而$3 \times 1$向量$\hat{\theta}$类似于"参数估计向量"。

式\eqref{eq:7.3.10}中定义的边界估计通过以下关系在线更新
\begin{equation}\label{eq:7.3.12}
\dot{\hat{\theta}} = \gamma S^T\|r\|
\end{equation}
其中$r$在式\eqref{eq:7.2.6}中定义，$S$在式\eqref{eq:7.3.10}中定义，$\gamma$是正的标量控制常数。为方便起见，我们还注意到，由于式\eqref{eq:7.3.4}中定义的$\delta_0$、$\delta_1$和$\delta_2$是常数，式\eqref{eq:7.3.12}可以写成
\begin{equation}\label{eq:7.3.13}
\dot{\tilde{\theta}} = -\gamma S^T\|r\|
\end{equation}
因为我们将定义$\theta$和$\hat{\theta}$之间的差为
\begin{equation}\label{eq:7.3.14}
\tilde{\theta} = \theta - \hat{\theta}
\end{equation}

我们现在将注意力转向分析式\eqref{eq:7.3.6}中给出的控制器的相应误差系统的稳定性。将控制器式\eqref{eq:7.3.6}代入机器人方程式\eqref{eq:7.3.1}得到误差系统
\begin{equation}\label{eq:7.3.15}
M(q)\dot{r} = -V_m(q,\dot{q})r - K_v r - v_R + w
\end{equation}
其中$w$在式\eqref{eq:7.3.2}中定义。

我们现在用Lyapunov-like函数分析式\eqref{eq:7.3.15}给出的误差系统的稳定性
\begin{equation}\label{eq:7.3.16}
V = \frac{1}{2}r^T M(q)r + \frac{1}{2}\tilde{\theta}^T \gamma^{-1}\tilde{\theta} + \frac{1}{2k_\varepsilon}\varepsilon^2
\end{equation}

对式\eqref{eq:7.3.16}关于时间求导得到
\begin{equation}\label{eq:7.3.17}
\dot{V} = \frac{1}{2}r^T\dot{M}(q)r + r^T M(q)\dot{r} + \tilde{\theta}^T\gamma^{-1}\dot{\tilde{\theta}} + \frac{1}{k_\varepsilon}\varepsilon\dot{\varepsilon}
\end{equation}
因为标量量可以转置。将式\eqref{eq:7.3.13}和\eqref{eq:7.3.15}代入式\eqref{eq:7.3.17}得到
\begin{equation}\label{eq:7.3.18}
\begin{aligned}
\dot{V} = {} & -r^T K_v r - (S\tilde{\theta})\|r\| + r^T(w - v_R) + \frac{1}{k_\varepsilon}\varepsilon\dot{\varepsilon} \\
& \quad {} + \frac{1}{2} r^T \bigl( \dot{M}(q) - 2V_m(q,\dot{q}) \bigr) r
\end{aligned}
\end{equation}

通过利用斜对称性质，很容易看出式\eqref{eq:7.3.18}中的第二行等于零。从式\eqref{eq:7.3.18}，我们可以使用式\eqref{eq:7.3.3}和\eqref{eq:7.3.11}以以下方式放置$\dot{V}$的上界：
\begin{equation}\label{eq:7.3.19}
\dot{V} \leq -r^T K_v r - (S\tilde{\theta})\|r\| + (S\theta)\|r\| - r^T v_R + \frac{1}{k_\varepsilon}\varepsilon\dot{\varepsilon}
\end{equation}

\addtocounter{equation}{1}% 略去书中一条中间恒等式，使式号自 (7.3.21) 起与原版一致

将式\eqref{eq:7.3.7}、\eqref{eq:7.3.8}、\eqref{eq:7.3.10}和\eqref{eq:7.3.14}代入式\eqref{eq:7.3.19}，我们得到
\begin{equation}\label{eq:7.3.21}
\dot{V} \leq -r^T K_v r - \varepsilon + S\hat{\theta}\|r\| - \frac{\|r\|^2 (S\hat{\theta})^2}{S\hat{\theta}\|r\| + \varepsilon}
\end{equation}

为式\eqref{eq:7.3.21}中的最后两项获得公共分母使我们能够将式\eqref{eq:7.3.21}写成
\begin{equation}\label{eq:7.3.22}
\dot{V} \leq -r^T K_v r - \varepsilon + \frac{\varepsilon S\hat{\theta}\|r\|}{S\hat{\theta}\|r\| + \varepsilon}
\end{equation}

由于式\eqref{eq:7.3.22}中最后两项之和总是小于零，我们可以在$\dot{V}$上放置新的上界：
\begin{equation}\label{eq:7.3.23}
\dot{V} \leq -r^T K_v r
\end{equation}

我们现在详细说明跟踪误差的稳定性类型。首先，注意从式\eqref{eq:7.3.23}，我们可以在$\dot{V}$上放置新的上界：
\begin{equation}\label{eq:7.3.24}
\dot{V} \leq -\lambda_{\min}\{K_v\}\|r\|^2
\end{equation}

如前一节对RCL所示，我们可以使用式\eqref{eq:7.3.24}来表明所有信号都是有界的，且$r \in L_2^n$（见第2章）。按照RCL稳定性分析，我们可以使用式\eqref{eq:7.2.6}来表明位置跟踪误差$e$与滤波跟踪误差$r$通过传递函数关系相关
\begin{equation}\label{eq:7.3.25}
e(s) = G(s)r(s)
\end{equation}
其中$s$是拉普拉斯变换变量，$G(s)$是严格适当的、渐近稳定的传递函数。因此，我们可以使用第2章中的定理2.4.7来说明
\begin{equation}\label{eq:7.3.26}
\lim_{t \to \infty} \|e\| = 0
\end{equation}

上面的结果告诉我们位置跟踪误差$e$是渐近稳定的。根据本节提出的理论开发，我们只能说明速度跟踪误差$\dot{e}$和边界估计$\hat{\theta}$是有界的。应该注意的是，在[Corless和Leitmann 1983]中，提出了更复杂的理论开发，证明速度跟踪误差是渐近稳定的。然而，为简洁起见，这些附加信息留给读者自行研究。

本节推导的自适应鲁棒控制器总结在表7.3.1中，并在图7.3.1中描绘。我们现在提出一个示例来说明如何使用表7.3.1为机器人机械臂设计自适应鲁棒控制器。

\textbf{表7.3.1}\quad 自适应鲁棒控制器

\begin{center}
\begin{tabular}{|l|l|}
\hline
\textbf{扭矩控制器} & $\tau = K_v r + \dfrac{r\hat{\rho}^2}{\hat{\rho}\|r\| + \varepsilon}$ \\
\hline
\textbf{滤波跟踪误差} & $r = e + \dot{e}$ \\
\hline
\textbf{$\varepsilon$ 动力学} & $\dot{\varepsilon} = -k_\varepsilon\varepsilon$ \\
\hline
\textbf{边界估计} & $\begin{aligned}[t]
\hat{\rho} &= S\hat{\theta},\\
S &= \begin{bmatrix} 1 & \|\mathbf{e}\| & \|\mathbf{e}\|^2 \end{bmatrix},\quad
\mathbf{e} = \begin{bmatrix} e \\ \dot{e} \end{bmatrix},\quad
\hat{\theta} = \begin{bmatrix} \hat{\delta}_0 \\ \hat{\delta}_1 \\ \hat{\delta}_2 \end{bmatrix}
\end{aligned}$ \\
\hline
\textbf{参数更新律} & $\dot{\hat{\theta}} = \gamma S^T\|r\|$ \\
\hline
\textbf{稳定性} & 位置跟踪误差$e$渐近稳定；边界估计$\hat{\theta}$有界；速度跟踪误差$\dot{e}$有界 \\
\hline
\textbf{注释} & 自适应鲁棒控制器无需修改即可补偿有界扰动 \\
\hline
\end{tabular}
\end{center}

\textbf{示例7.3--1：二连杆机械臂的自适应鲁棒控制器}

我们希望为图6.2.1中给出的二连杆机械臂设计和仿真表7.3.1中给出的自适应鲁棒控制器。（该机器人臂的动力学在第3章中给出。）为了建模摩擦与扰动，分别在关节扭矩中加入以下项：
\begin{equation}
\begin{aligned}
&2\dot{q}_1 + 0.5\operatorname{sgn}(\dot{q}_1) + 0.2\sin(3t),\\
&2\dot{q}_2 + 0.5\operatorname{sgn}(\dot{q}_2) + 0.2\sin(3t),
\end{aligned}
\end{equation}
其中第一式对应$\tau_1$，第二式对应$\tau_2$。

我们现在可以使用表7.3.1来制定自适应鲁棒控制器为
\begin{align}
\tau_1 &= k_v r_1 + \frac{r_1\hat{\rho}^2}{\hat{\rho}\|r\| + \varepsilon}, \\
\tau_2 &= k_v r_2 + \frac{r_2\hat{\rho}^2}{\hat{\rho}\|r\| + \varepsilon},
\end{align}
其中$K_v = k_v I$（$I$为$2\times 2$单位矩阵），$r_1 = e_1 + \dot{e}_1$，$r_2 = e_2 + \dot{e}_2$，$\dot{\varepsilon} = -k_\varepsilon\varepsilon$，且
\begin{equation}
\|r\| = \sqrt{r_1^2 + r_2^2}.
\end{equation}

在上面控制扭矩的表达式中，边界函数$\hat{\rho}$写为
\begin{equation}
\hat{\rho} = S\hat{\theta}
= \begin{bmatrix} 1 & \|e\| & \|e\|^2 \end{bmatrix}
\begin{bmatrix} \hat{\delta}_0 \\ \hat{\delta}_1 \\ \hat{\delta}_2 \end{bmatrix},
\end{equation}
其中
\begin{equation}
\|e\| = \sqrt{e_1^2 + e_2^2 + \dot{e}_1^2 + \dot{e}_2^2},
\end{equation}
而$\hat{\delta}_0$、$\hat{\delta}_1$和$\hat{\delta}_2$为自适应估计。从表7.3.1，相关边界估计按以下方式更新：
\begin{equation}
\dot{\hat{\delta}}_0 = \gamma\|r\|, \quad \dot{\hat{\delta}}_1 = \gamma\|e\|\|r\|, \quad \dot{\hat{\delta}}_2 = \gamma\|e\|^2\|r\|,
\end{equation}
其中$\gamma$为正的标量自适应增益。

对于$m_1=0.8$ kg，$m_2=2.3$ kg，每根连杆长度为1 m，自适应鲁棒控制器在以下条件下进行了仿真：控制参数、初始条件和期望轨迹由
\begin{equation*}
k_v = 50, \quad \gamma = 10, \quad k_\varepsilon = 0.001, \quad \varepsilon(0) = 1
\end{equation*}
和
\begin{equation*}
q_{d1} = \sin t, \quad q_{d2} = \cos t, \quad \hat{\delta}_0(0) = 0, \quad \hat{\delta}_1(0) = 0, \quad \hat{\delta}_2(0) = 0
\end{equation*}
给出

跟踪误差与边界估计在图7.3.2中描绘。如图所示，位置和速度跟踪误差都是渐近稳定的，边界估计保持有界。应该注意的是，从本节给出的理论开发，我们只保证位置跟踪误差是渐近稳定的，而所有其他信号保持有界。

\section{执行器动力学补偿}

在本书中，我们讨论了在"扭矩输入级别"设计的控制器。也就是说，与执行器相关的任何动力学都被忽略了。提出这一点的原因不是要贬低先前讨论的控制开发，因为这项研究涉及解决一个非常困难的问题，即在不确定性存在的情况下高度非线性系统的全局跟踪控制。我们提到以前方法中的这一不足是为了强调这样一个事实：根据许多机器人控制研究人员的观点，现在是时候开始在控制综合中包括执行器动力学的效应了。最近，一些研究人员假设执行器的有害效应正在阻止机器人机械臂的高速运动/力控制[Eppinger和Seering 1987]。

在本节中，我们说明了如何使用系统方法来补偿电气效应和关节柔性形式的执行器动力学。使用这种方法和精确模型知识的假设，开发了产生连杆跟踪误差全局渐近稳定性结果的控制器。尽管我们假设精确知道模型，但重要的是要认识到，在某些情况下，可能可以制定自适应和鲁棒非线性跟踪控制器来补偿"不确定性"。在存在执行器动力学的情况下补偿不确定系统目前正在研究中[Ghorbel和Spong 1990]。

\subsection{电气动力学}

在本小节中，我们说明了如何合成"校正"控制器[Kokotovic等1986]，以确保尽管电机会增加整体系统动力学的电气动力学，连杆跟踪也是渐近的。"校正控制器"一词用于强调控制器校正电气动力学。本小节中研究的机器人类将被称为刚连杆电气驱动（RLED）机器人。为简单起见，我们假设执行器是直流（dc）电机；然而，通过一些修改，以下分析可用于更复杂的电机，如开关磁阻电机[Taylor 1989]。

RLED机器人的模型[Tarn等1991]采用
\begin{equation}\label{eq:7.4.1}
\bar{M}(q)\ddot{q} + N(q,\dot{q}) = K_T I
\end{equation}
\begin{equation}\label{eq:7.4.2}
L_a\dot{I} + R(I,\dot{q}) = u_E
\end{equation}
其中
\begin{equation}\label{eq:7.4.3}
\bar{M}(q) = M(q) + J
\end{equation}
$M(q)$是$n \times n$连杆惯性矩阵，$N(q,\dot{q})$是包含向心力、科里奥利力、重力、阻尼和摩擦项的$n \times 1$向量，$J$是用于表示执行器惯性的$n \times n$常数、对角、正定矩阵，$I(t)$是用于表示每个执行器中电流的$n \times 1$向量，$K_T$是用于表示扭矩和电流之间转换的常数对角$n \times n$矩阵，$L_a$是用于表示电气电感的$n \times n$常数正定对角矩阵，$R(q,\dot{q})$是用于表示电气电阻和电机反电动势的$n \times 1$向量，$u_E(t)$是用于表示输入电机电压的$n \times 1$控制向量。

在本书中，重点放在利用机器人机械臂的物理特性来帮助我们进行稳定性分析。遵循这一传统，我们注意到式\eqref{eq:7.4.3}中定义的复合惯性矩阵$\bar{M}(q)$是对称的、正定的，并且作为$q$的函数一致有界；因此，我们可以对任何$n \times 1$向量$x$说明
\begin{equation}\label{eq:7.4.4}
m_1 \|x\|^2 = \lambda_{\min}\{\bar{M}(q)\}\|x\|^2 \leq x^T \bar{M}(q) x
\end{equation}
其中$m_1$是正的标量常数，取决于特定机器人的质量特性（见第3章）。从式\eqref{eq:7.4.4}，也可以建立
\begin{equation}\label{eq:7.4.5}
x^T \bar{M}^{-1}(q) x \leq \lambda_{\max}\{\bar{M}^{-1}(q)\}\|x\|^2 = \frac{1}{m_1}\|x\|^2 = \|\bar{M}^{-1}(q)\|_{i2}\|x\|^2
\end{equation}
其中$\|\cdot\|_{i2}$用于表示诱导2-范数（见第2章）。

如多次讨论的那样，我们对连杆跟踪误差的性能感兴趣。为避免混淆，我们重申跟踪误差定义为
\begin{equation}\label{eq:7.4.6}
e = q_d - q
\end{equation}
其中$q_d$表示期望连杆轨迹。我们将假设$q_d$及其第一、第二和第三导数都是作为时间函数有界的。我们还假设式\eqref{eq:7.4.1}左侧连杆动力学的第一导数存在。对期望轨迹和连杆动力学的这些"平滑性"假设确保稍后开发的控制器保持有界。

控制目标将是尽管有电气动力学，也要获得渐近连杆跟踪。为了实现这一目标，我们首先用式\eqref{eq:7.4.6}给出的跟踪误差将式\eqref{eq:7.4.1}重写为
\begin{equation}\label{eq:7.4.7}
\bar{M}(q)\ddot{q}_d - \bar{M}(q)\ddot{e} + N(q,\dot{q}) = K_T I
\end{equation}
式\eqref{eq:7.4.7}给出的误差系统也可以用状态空间形式写成
\begin{equation}\label{eq:7.4.8}
\dot{\mathbf{e}} = A_o \mathbf{e} + B\left[\ddot{q}_d + \bar{M}^{-1}(q)N(q,\dot{q}) - \bar{M}^{-1}(q)K_T I\right]
\end{equation}
其中
\begin{equation*}
\mathbf{e} = \begin{bmatrix} e \\ \dot{e} \end{bmatrix}, \quad
A_o = \begin{bmatrix} O_{n \times n} & I_{n \times n} \\ O_{n \times n} & O_{n \times n} \end{bmatrix}, \quad
B = \begin{bmatrix} O_{n \times n} \\ I_{n \times n} \end{bmatrix}
\end{equation*}
$O_{n \times n}$是$n \times n$零矩阵，$I_{n \times n}$是$n \times n$单位矩阵。

如人们可以清楚地看到，式\eqref{eq:7.4.8}中没有控制输入；因此，我们将在式\eqref{eq:7.4.8}的右侧添加和减去项$B\bar{M}^{-1}(q)u_L$得到
\begin{equation}\label{eq:7.4.9}
\dot{\mathbf{e}} = A_o \mathbf{e} + B\left[\ddot{q}_d + \bar{M}^{-1}(q)N(q,\dot{q}) - \bar{M}^{-1}(q)u_L\right] + B\left[\bar{M}^{-1}(q)(u_L - K_T I)\right]
\end{equation}
其中$u_L$是表示"虚构"$n \times 1$控制输入的$n \times 1$向量。事实证明，若电气动力学不存在，则$u_L$可取为计算扭矩控制器并保证连杆跟踪误差渐近稳定。正如我们稍后将看到的，虚构控制器$u_L$实际上嵌入在整体控制策略中，该策略设计在电压控制输入$u_E$处。

继续误差系统开发，我们为RLED机器人定义$u_L$为计算扭矩控制器
\begin{equation}\label{eq:7.4.10}
u_L = \bar{M}(q)(\ddot{q}_d + K_{Lv}\dot{e} + K_{Lp}e) + N(q,\dot{q})
\end{equation}
其中$K_{Lv}$和$K_{Lp}$定义为$n \times n$正定对角矩阵。仅将式\eqref{eq:7.4.10}的第一个$u_L$项代入式\eqref{eq:7.4.9}得到连杆跟踪误差系统
\begin{equation}\label{eq:7.4.11}
\dot{\mathbf{e}} = A_L \mathbf{e} + B\bar{M}^{-1}(q)\eta_E
\end{equation}
其中
\begin{equation}\label{eq:7.4.12}
A_L = \begin{bmatrix} O_{n \times n} & I_{n \times n} \\ -K_{Lp} & -K_{Lv} \end{bmatrix}, \quad \eta_E = u_L - K_T I
\end{equation}

关于式\eqref{eq:7.4.11}给出的连杆跟踪误差系统，如果$\eta_E$可以保证在所有时间为零，我们可以很容易地表明跟踪误差是渐近稳定的，因为式\eqref{eq:7.4.12}中定义的$A_L$具有稳定的特征值。因此，可以将控制目标视为迫使"扰动"$\eta_E$趋近于零。

为了设计$\eta_E$的控制律，我们必须首先建立其动态特性。由式\eqref{eq:7.4.12}中$\eta_E = u_L - K_T I$，有
\begin{equation}\label{eq:7.4.13}
\dot{\eta}_E = \dot{u}_L - K_T \dot{I}
\end{equation}

为了获得$\eta_E$的动态特性，我们将式\eqref{eq:7.4.2}中的$\dot{I}$代入式\eqref{eq:7.4.13}得到
\begin{equation}\label{eq:7.4.14}
\dot{\eta}_E = \dot{u}_L - K_T L_a^{-1}(u_E - R(I,\dot{q}))
\end{equation}

我们现在可以使用式\eqref{eq:7.4.14}设计输入$u_E$的控制律以迫使$\eta_E$趋近于零。$\eta_E$应该趋近于零的事实推动了校正控制律
\begin{equation}\label{eq:7.4.15}
u_E = L_a K_T^{-1}(\dot{u}_L + K_{Ep}\eta_E) + R(I,\dot{q})
\end{equation}
其中$K_{Ep}$定义为$n \times n$正定对角矩阵。将式\eqref{eq:7.4.15}代入式\eqref{eq:7.4.14}得到
\begin{equation}\label{eq:7.4.16}
\dot{\eta}_E = -K_{Ep}\eta_E
\end{equation}

式\eqref{eq:7.4.11}和\eqref{eq:7.4.16}给出的动力学方程可以被认为是代表整体闭环动力学的两个相互连接的子系统。如人们所期望的，确定整体闭环系统的稳定性类型将是可取的。为了确定式\eqref{eq:7.4.11}和\eqref{eq:7.4.16}的稳定性类型，我们将利用Lyapunov函数
\begin{equation}\label{eq:7.4.17}
V = \mathbf{e}^T P_L \mathbf{e} + \frac{1}{2}\eta_E^T \eta_E
\end{equation}
其中
\begin{equation*}
P_L = \frac{1}{2}\begin{bmatrix} K_{Lp} + \frac{1}{2}K_{Lv} & \frac{1}{2}I_{n \times n} \\ \frac{1}{2}I_{n \times n} & I_{n \times n} \end{bmatrix}
\end{equation*}

如果满足
\begin{equation}\label{eq:7.4.18}
\lambda_{\min}\{K_{Lv}\} > 1
\end{equation}
给出的充分条件，则由Gerschgorin定理（见第2章），很明显$P_L$是正定矩阵，因此$V$是Lyapunov函数。式\eqref{eq:7.4.18}给出的条件简单地意味着最小速度控制器增益应大于1。

对式\eqref{eq:7.4.17}关于时间求导得到
\begin{equation}\label{eq:7.4.19}
\dot{V} = \mathbf{e}^T P_L \dot{\mathbf{e}} + \dot{\mathbf{e}}^T P_L \mathbf{e} + \frac{1}{2}\dot{\eta}_E^T \eta_E + \frac{1}{2}\eta_E^T \dot{\eta}_E
\end{equation}
将式\eqref{eq:7.4.11}和\eqref{eq:7.4.16}代入式\eqref{eq:7.4.19}得到
\begin{equation}\label{eq:7.4.20}
\dot{V} = -\mathbf{e}^T Q_L \mathbf{e} - \eta_E^T K_{Ep}\eta_E + 2\mathbf{e}^T P_L B \bar{M}^{-1}(q)\eta_E
\end{equation}
其中
\begin{equation}\label{eq:7.4.21}
Q_L = -(A_L^T P_L + P_L A_L) = \begin{bmatrix} \frac{1}{2}K_{Lp} & O_{n \times n} \\ O_{n \times n} & K_{Lv} - \frac{1}{2}I_{n \times n} \end{bmatrix}
\end{equation}

注意，如果式\eqref{eq:7.4.18}给出的充分条件成立，很明显式\eqref{eq:7.4.20}中的矩阵$Q_L$是正定的。利用$Q_L$是正定的事实使我们能够在式\eqref{eq:7.4.20}中给出的$\dot{V}$上放置上界。这个上界由
\begin{equation}\label{eq:7.4.22}
\dot{V} \leq -\lambda_{\min}\{Q_L\}\|\mathbf{e}\|^2 - \lambda_{\min}\{K_{Ep}\}\|\eta_E\|^2 + 2\beta_o \|\mathbf{e}\|\|\eta_E\|
\end{equation}
给出，其中
\begin{equation*}
\beta_o = \|P_L B \bar{M}^{-1}(q)\|_{i2} = \frac{1}{2}\left\|\begin{bmatrix} \frac{1}{2}\bar{M}^{-1}(q) \\ \bar{M}^{-1}(q) \end{bmatrix}\right\|_{i2}
\end{equation*}
而$m_1$及$\|\bar{M}^{-1}(q)\|_{i2}=1/m_1$见式\eqref{eq:7.4.4}与式\eqref{eq:7.4.5}。

为了确定渐近稳定性的控制器增益的充分条件，我们将式\eqref{eq:7.4.22}写成矩阵形式
\begin{equation}\label{eq:7.4.23}
\dot{V} \leq -\begin{bmatrix} \|\mathbf{e}\| & \|\eta_E\| \end{bmatrix} Q_o \begin{bmatrix} \|\mathbf{e}\| \\ \|\eta_E\| \end{bmatrix}
\end{equation}
其中
\begin{equation*}
Q_o = \begin{bmatrix} \lambda_{\min}\{Q_L\} & -\beta_o \\ -\beta_o & \lambda_{\min}\{K_{Ep}\} \end{bmatrix}, \quad
x_o = \begin{bmatrix} \|\mathbf{e}\| \\ \|\eta_E\| \end{bmatrix}
\end{equation*}
（原文献亦将$\|\mathbf{e}\|$写作$\|e\|$。）

由Gerschgorin定理，如果满足
\begin{equation}\label{eq:7.4.24}
\min\left\{ \lambda_{\min}\{Q_L\},\, \lambda_{\min}\{K_{Ep}\} \right\} > \beta_o
\end{equation}
给出的充分条件，则式\eqref{eq:7.4.23}中定义的矩阵$Q_o$将为正定。因此，如果控制器增益满足式\eqref{eq:7.4.18}和\eqref{eq:7.4.24}给出的条件，我们可以使用标准Lyapunov稳定性论证（见第2章）来说明$x_o$为渐近稳定，从而$\|\mathbf{e}\|$、$e$与$\dot{e}$均为渐近稳定。很容易表明，若满足
\begin{equation}\label{eq:7.4.25}
\min\left\{ \lambda_{\min}\{K_{Lp}\},\, \lambda_{\min}\{K_{Lv}\},\, \lambda_{\min}\{K_{Ep}\} \right\} > \frac{2}{m_1} + 1
\end{equation}
则式\eqref{eq:7.4.18}与式\eqref{eq:7.4.24}给出的增益条件均被满足。

应该注意的是，式\eqref{eq:7.4.15}给出的控制表面上似乎需要测量位置、速度、加速度及电流等信息；然而，由于我们假设精确知道式\eqref{eq:7.4.1}和\eqref{eq:7.4.2}给出的动力学模型，我们可以消除对加速度$\ddot{q}$的直接测量需求。也就是说，通过对式\eqref{eq:7.4.10}关于时间求导，可将$\dot{u}_L$写成
\begin{equation}\label{eq:7.4.26}
\begin{aligned}
\dot{u}_L &= \dot{\bar{M}}(q)(\ddot{q}_d + K_{Lv}\dot{e} + K_{Lp}e) + \dot{N}(q,\dot{q}) \\
&\quad {} + \bar{M}(q)\left[ \frac{d}{dt}\ddot{q}_d + K_{Lv}(\ddot{q}_d - \ddot{q}) + K_{Lp}\dot{e} \right]
\end{aligned}
\end{equation}
其中$\ddot{q}$从式\eqref{eq:7.4.1}中找到为
\begin{equation*}
\ddot{q} = \bar{M}^{-1}(q)[K_T I - N(q,\dot{q})]
\end{equation*}

在式\eqref{eq:7.4.26}中代入$\ddot{q}$后，$\dot{u}_L$将仅取决于$q$、$\dot{q}$和$I$的测量。在控制输入$u_E$处实现的实际控制可以通过进行适当的代入找到。也就是说，式\eqref{eq:7.4.15}给出的校正控制可以写成
\begin{equation}\label{eq:7.4.27}
\begin{aligned}
u_E &= L_a K_T^{-1}\dot{u}_L + R(I,\dot{q}) \\
&\quad {} + L_a K_T^{-1}K_{Ep}\bigl[ \bar{M}(q)(\ddot{q}_d + K_{Lv}\dot{e} + K_{Lp}e) + N(q,\dot{q}) - K_T I \bigr]
\end{aligned}
\end{equation}
其中$\dot{u}_L$由式\eqref{eq:7.4.26}给出（方括号内即为$\eta_E=u_L-K_TI$）。

在仔细检查式\eqref{eq:7.4.26}中给出的$\dot{u}_L$的函数依赖性后，现在很明显为什么我们假设期望轨迹和连杆动力学足够平滑。具体而言，从式\eqref{eq:7.4.26}我们可以看到，校正控制器需要期望轨迹的第一、第二和第三时间导数有界，同时需要连杆动力学的第一导数存在。对期望轨迹和连杆动力学的这些假设确保控制输入保持有界。

上面推导的校正控制器总结在表7.4.1中，并在图7.4.1中描绘。

\textbf{表7.4.1}\quad RLED校正控制器

\begin{center}
\small
\begin{tabular}{|l|p{11.2cm}|}
\hline
\textbf{电压控制器} & \begin{minipage}[t]{\linewidth}\raggedright\setlength{\lineskip}{6pt}
$\begin{aligned}[t] u_E &= L_a K_T^{-1}\dot{u}_L + R(I,\dot{q}) \\ &\quad {} + L_a K_T^{-1}K_{Ep}\bigl[\bar{M}(q)(\ddot{q}_d + K_{Lv}\dot{e} + K_{Lp}e) \\ &\qquad {} + N(q,\dot{q}) - K_T I\bigr] \end{aligned}$\par
$\begin{aligned}[t] \dot{u}_L &= \dot{\bar{M}}(q)(\ddot{q}_d + K_{Lv}\dot{e} + K_{Lp}e) + \dot{N}(q,\dot{q}) \\ &\quad {} + \bar{M}(q)\left[\frac{\mathrm{d}}{\mathrm{d}t}\ddot{q}_d + K_{Lv}(\ddot{q}_d - \ddot{q}) + K_{Lp}\dot{e}\right] \end{aligned}$\par
$\ddot{q} = \bar{M}^{-1}(q)\bigl[K_T I - N(q,\dot{q})\bigr]$
\end{minipage} \\
\hline
\textbf{稳定性} & 跟踪误差$e$与$\dot{e}$渐近稳定。 \\
\hline
\textbf{注释} & 控制器需要系统动力学的精确知识；增益矩阵$K_{Lv}$、$K_{Lp}$与$K_{Ep}$须足够大；期望轨迹须足够光滑。 \\
\hline
\end{tabular}
\end{center}

其中$\dot{u}_L$亦可通过对$u_L=\bar{M}(q)(\ddot{q}_d + K_{Lv}\dot{e} + K_{Lp}e)+N(q,\dot{q})$求导并将式\eqref{eq:7.4.1}代入$\ddot{q}$得到，与表中第二式等价。

我们现在提出一个示例来说明如何使用表7.4.1为RLED机器人设计校正控制器。

\textbf{示例7.4--1：单连杆RLED机械臂的校正控制器}

我们希望使用表7.4.1为图7.4.2中给出的单连杆电机驱动机器人臂设计和仿真校正控制器。系统的动力学采用
\begin{align}
(mL^2 + J)\ddot{q} + mLg\sin q + f_d \dot{q} &= K_T I, \tag{1} \\
L_a \dot{I} + RI + k_b \dot{q} &= u_E. \tag{2}
\end{align}
其中$m=1$ kg，$K_T=2$ N/A，$k_b=0.3$ V-S，$f_d=3$ kg-m/s，$L=1$ m，$L_a=0.1$ H，$R=5\ \Omega$，$g$是重力系数，$J=0.2$ kg-m$^2$，$I$是电机电流，$u_E$是电机输入电压。

假设式(1)和(2)给出的模型精确知道，我们可以使用表7.4.1来制定校正控制器为
\begin{align}
u_E &= L_a K_T^{-1}\dot{u}_L + RI + k_b\dot{q} \notag \\
&\quad {} + L_a K_T^{-1}K_{Ep}\bigl[(mL^2 + J)(\ddot{q}_d + K_{Lv}\dot{e} + K_{Lp}e) + mLg\sin q + f_d\dot{q} - K_T I\bigr], \tag{3} \\
\dot{u}_L &= mLg\dot{q}\cos q + f_d\ddot{q} + (mL^2 + J)\left[\frac{d}{dt}\ddot{q}_d + K_{Lv}(\ddot{q}_d - \ddot{q}) + K_{Lp}\dot{e}\right], \tag{4}
\end{align}
由式(1)，$\ddot{q}$可写为
\begin{equation}
\ddot{q} = (mL^2 + J)^{-1}\bigl[K_T I - mLg\sin q - f_d\dot{q}\bigr]. \tag{5}
\end{equation}
将式(5)代入式(4)，$\dot{u}_L$即可仅用可测量与模型参数表示。

校正控制器在以下条件下进行了仿真：控制参数、初始条件和期望轨迹由
\begin{equation*}
K_{Lp} = K_{Lv} = K_{Ep} = 5, \quad q(0) = \dot{q}(0) = I(0) = 0
\end{equation*}
和
\begin{equation*}
q_d = \sin t
\end{equation*}
给出

跟踪误差和控制电压在图7.4.3中描绘。如图所示，跟踪误差是渐近稳定的。

\subsection{关节柔性}

在本小节中，我们说明了如何合成"校正"控制器，以确保尽管驱动或齿轮会增加整体系统动力学的关节柔性，连杆跟踪也是渐近的。"校正控制器"一词用于强调控制器校正用于表示关节柔性效应的动力学。本小节中研究的机器人类将被称为刚连杆柔性关节（RLFJ）机器人。

RLFJ机器人的模型[Spong 1987]取为
\begin{equation}\label{eq:7.4.28}
M(q)\ddot{q} + \mathbf{N}(q,\dot{q}) = K q_m
\end{equation}
\begin{equation}\label{eq:7.4.29}
J\ddot{q}_m + B_m(q_m,\dot{q}_m,q) = u_F,
\end{equation}
其中
\begin{equation}\label{eq:7.4.30}
\mathbf{N}(q,\dot{q}) = N(q,\dot{q}) + K q,
\end{equation}
$q_m(t)$为电机位移的$n \times 1$向量，$K$为常数、对角、正定关节刚度矩阵，$M(q)$为连杆惯性矩阵，$N(q,\dot{q})$表示刚体机器人的非线性项（科氏力、离心力与重力等），$J$为电机惯性矩阵，$B_m(q_m,\dot{q}_m,q)$汇集电机侧阻尼及柔性耦合所产生的扭矩项，$u_F$为$n \times 1$关节输入扭矩；其余记号与前一子节一致。将式\eqref{eq:7.4.30}代入式\eqref{eq:7.4.28}可得与文献中常见的$K(q_m-q)$形式等价的连杆加速度表达式。关于$M(q)$，我们从第3章注意到，对于任何$n \times 1$向量$x$
\begin{equation}\label{eq:7.4.31}
x^{\mathsf T} M^{-1}(q) x \leq \lambda_{\max}\{M^{-1}(q)\}\|x\|^2 = \frac{1}{m_1}\|x\|^2 = \|M^{-1}(q)\|_{i2}\|x\|^2
\end{equation}
其中$m_1$为正的标量常数（可取为$\lambda_{\min}\{M(q)\}$的相应界限，与下文Lyapunov分析中记号一致）。

如前一子节，我们对式\eqref{eq:7.4.6}中定义的连杆跟踪误差的性能感兴趣。对于RLFJ机器人的控制，我们将假设$q_d$及其第一、第二、第三和第四导数都是作为时间函数有界的。我们还假设式\eqref{eq:7.4.28}左侧连杆动力学的第一和第二导数存在。对期望轨迹和连杆动力学的这些"平滑性"假设确保稍后开发的控制器保持有界。

遵循前面各节给出的相同分析开发，我们用式\eqref{eq:7.4.6}给出的跟踪误差将式\eqref{eq:7.4.28}写成状态空间形式
\begin{equation}\label{eq:7.4.32}
\dot{\mathbf{e}} = A_o \mathbf{e} + B\left[\ddot{q}_d + M^{-1}(q)\mathbf{N}(q,\dot{q}) - M^{-1}(q)K q_m\right]
\end{equation}
其中$\mathbf{e}$、$A_o$与$B$由式\eqref{eq:7.4.8}给出。同样，由于式\eqref{eq:7.4.32}中没有虚构控制输入；我们在式\eqref{eq:7.4.32}的右侧添加和减去项$B\,M^{-1}(q)\,u_L$得到
\begin{equation}\label{eq:7.4.33}
\begin{aligned}
\dot{\mathbf{e}} &= A_o \mathbf{e} + B\left[\ddot{q}_d + M^{-1}(q)\mathbf{N}(q,\dot{q}) - M^{-1}(q)u_L\right] \\
&\quad + B\left[M^{-1}(q)(u_L - K q_m)\right]
\end{aligned}
\end{equation}
其中$u_L$再次用于表示虚构的$n \times 1$控制输入。如前所述，虚构控制器$u_L$将嵌入在整体控制策略中，该策略设计在控制输入$u_F$处。

继续误差系统开发，我们为RLFJ机器人定义$u_L$为计算扭矩控制器
\begin{equation}\label{eq:7.4.34}
u_L = M(q)(\ddot{q}_d + K_{Lv}\dot{e} + K_{Lp}e) + \mathbf{N}(q,\dot{q})
\end{equation}
其中$K_{Lv}$和$K_{Lp}$如式\eqref{eq:7.4.10}中定义。将式\eqref{eq:7.4.34}代入式\eqref{eq:7.4.33}得到连杆跟踪误差系统
\begin{equation}\label{eq:7.4.35}
\dot{\mathbf{e}} = A_L \mathbf{e} + B M^{-1}(q) C \boldsymbol{\eta}_F
\end{equation}
其中$A_L$如式\eqref{eq:7.4.12}中定义，
\begin{equation}\label{eq:7.4.36}
C = \begin{bmatrix} I_{n \times n} & O_{n \times n} \end{bmatrix}, \quad \boldsymbol{\eta}_F = \begin{bmatrix} \eta_F \\ \dot{\eta}_F \end{bmatrix} = \begin{bmatrix} u_L - K q_m \\ \dot{u}_L - K \dot{q}_m \end{bmatrix}
\end{equation}

以式\eqref{eq:7.4.36}中给出的方式定义$\boldsymbol{\eta}_F$（及其分量$\eta_F$、$\dot{\eta}_F$）的原因在于式\eqref{eq:7.4.29}中给出的动力学是二阶动力学。也就是说，由于执行器动力学是二阶的，我们迫使$\boldsymbol{\eta}_F$及其导数$\dot{\boldsymbol{\eta}}_F$为零以确保连杆跟踪误差$e$趋近于零。

为了设计$\boldsymbol{\eta}_F$的控制律，我们必须首先建立其动态特性。由式\eqref{eq:7.4.36}，有
\begin{equation}\label{eq:7.4.37}
\dot{\boldsymbol{\eta}}_F = \begin{bmatrix} \dot{\eta}_F \\ \ddot{\eta}_F \end{bmatrix} = \begin{bmatrix} \dot{u}_L - K \dot{q}_m \\ \ddot{u}_L - K \ddot{q}_m \end{bmatrix}
\end{equation}

为了获得$\boldsymbol{\eta}_F$的动态特性，我们将式\eqref{eq:7.4.29}中的$\ddot{q}_m$代入式\eqref{eq:7.4.37}得到
\begin{equation}\label{eq:7.4.38}
\dot{\boldsymbol{\eta}}_F = \begin{bmatrix} \dot{u}_L - K \dot{q}_m \\ \ddot{u}_L - K J^{-1}\bigl(u_F - B_m(q_m,\dot{q}_m,q)\bigr) \end{bmatrix}
\end{equation}

我们现在可以使用式\eqref{eq:7.4.38}设计输入$u_F$的控制律以迫使$\boldsymbol{\eta}_F$趋近于零。$\boldsymbol{\eta}_F$应该趋近于零的事实推动了控制律
\begin{equation}\label{eq:7.4.39}
u_F = J K^{-1}\bigl(\ddot{u}_L + K_{Fv}\dot{\eta}_F + K_{Fp}\eta_F\bigr) + B_m(q_m,\dot{q}_m,q)
\end{equation}
其中$K_{Fv}$和$K_{Fp}$定义为$n \times n$正定对角矩阵（式中的$\eta_F$、$\dot{\eta}_F$为式\eqref{eq:7.4.36}所定义之分量）。将式\eqref{eq:7.4.39}代入式\eqref{eq:7.4.38}得到
\begin{equation}\label{eq:7.4.40}
\dot{\boldsymbol{\eta}}_F = A_F \boldsymbol{\eta}_F
\end{equation}
\begin{equation*}
A_F = \begin{bmatrix} \boldsymbol{0}_{n \times n} & I_{n \times n} \\ -K_{Fp} & -K_{Fv} \end{bmatrix}
\end{equation*}

式\eqref{eq:7.4.35}和式\eqref{eq:7.4.40}给出的动力学方程可以被认为是代表整体闭环动力学的两个相互连接的子系统。为了确定闭环动力学的稳定性类型，我们将利用Lyapunov函数
\begin{equation}\label{eq:7.4.41}
V = \mathbf{e}^{\mathsf T} P_L \mathbf{e} + \boldsymbol{\eta}_F^{\mathsf T} P_F \boldsymbol{\eta}_F
\end{equation}
其中$P_L$在式\eqref{eq:7.4.17}中定义且
\begin{equation*}
P_F = \frac{1}{2} \begin{bmatrix} K_{Fp} + \frac{1}{2}K_{Fv} & \frac{1}{2}I_{n \times n} \\ \frac{1}{2}I_{n \times n} & I_{n \times n} \end{bmatrix}
\end{equation*}

如果满足
\begin{equation}\label{eq:7.4.42}
\min\left\{\lambda_{\min}\{K_{Lv}\},\lambda_{\min}\{K_{Fv}\}\right\} > 1
\end{equation}
给出的充分条件，则由Gerschgorin定理很明显式\eqref{eq:7.4.41}中的矩阵$P_L$和$P_F$是正定矩阵，因此$V$是Lyapunov函数。

对式\eqref{eq:7.4.41}关于时间求导得到
\begin{equation}\label{eq:7.4.43}
\dot{V} = \mathbf{e}^{\mathsf T} P_L \dot{\mathbf{e}} + \dot{\mathbf{e}}^{\mathsf T} P_L \mathbf{e} + \boldsymbol{\eta}_F^{\mathsf T} P_F \dot{\boldsymbol{\eta}}_F + \dot{\boldsymbol{\eta}}_F^{\mathsf T} P_F \boldsymbol{\eta}_F
\end{equation}
将式\eqref{eq:7.4.35}和式\eqref{eq:7.4.40}代入式\eqref{eq:7.4.43}得到
\begin{equation}\label{eq:7.4.44}
\dot{V} = -\mathbf{e}^{\mathsf T} Q_L \mathbf{e} - \boldsymbol{\eta}_F^{\mathsf T} Q_F \boldsymbol{\eta}_F + 2\, \mathbf{e}^{\mathsf T} P_L B M^{-1}(q) C \boldsymbol{\eta}_F
\end{equation}
其中$Q_L$在式\eqref{eq:7.4.21}中定义且
\begin{equation}\label{eq:7.4.45}
Q_F = -(A_F^{\mathsf T} P_F + P_F A_F) = \begin{bmatrix} \frac{1}{2}K_{Fp} & O_{n \times n} \\ O_{n \times n} & K_{Fv} - \frac{1}{2}I_{n \times n} \end{bmatrix}
\end{equation}

注意，如果式\eqref{eq:7.4.42}给出的充分条件成立，很明显式\eqref{eq:7.4.44}中的矩阵$Q_L$和$Q_F$是正定矩阵。利用$Q_L$和$Q_F$是正定的事实使我们能够在式\eqref{eq:7.4.44}中给出的$\dot{V}$上放置上界。这个上界由
\begin{equation}\label{eq:7.4.46}
\dot{V} \leq -\lambda_{\min}\{Q_L\}\|\mathbf{e}\|^2 - \lambda_{\min}\{Q_F\}\|\boldsymbol{\eta}_F\|^2 + 2\, \beta_1 \|\mathbf{e}\| \|\boldsymbol{\eta}_F\|
\end{equation}
给出，其中
\begin{equation*}
\beta_1 = \frac{1}{m_1} > \left\|P_L B M^{-1}(q) C\right\|_{i2} = \frac{1}{2}\left\|\begin{bmatrix} \frac{1}{2}M^{-1}(q) & O_{n \times n} \\ M^{-1}(q) & O_{n \times n} \end{bmatrix}\right\|_{i2}
\end{equation*}
其中$m_1$在式\eqref{eq:7.4.31}中定义。

为了确定渐近稳定性的控制器增益的充分条件，我们将式\eqref{eq:7.4.46}写成矩阵形式
\begin{equation}\label{eq:7.4.47}
\dot{V} \leq -x_1^{\mathsf T} Q_1 x_1
\end{equation}
其中
\begin{equation*}
Q_1 = \begin{bmatrix} \lambda_{\min}\{Q_L\} & -\beta_1 \\ -\beta_1 & \lambda_{\min}\{Q_F\} \end{bmatrix}, \quad x_1 = \begin{bmatrix} \|e\| \\ \|\eta_F\| \end{bmatrix}
\end{equation*}

由Gerschgorin定理，如果满足
\begin{equation}\label{eq:7.4.48}
\min\left\{\lambda_{\min}\{Q_L\},\lambda_{\min}\{Q_F\}\right\} > \beta_1
\end{equation}
给出的充分条件，则式\eqref{eq:7.4.47}中定义的矩阵$Q_1$将为正定。因此，如果控制器增益满足式\eqref{eq:7.4.42}和式\eqref{eq:7.4.48}给出的条件，我们可以使用标准Lyapunov稳定性论证来说明式\eqref{eq:7.4.47}中定义的向量$x_1$，从而$\|e\|$、$e$及$\dot{e}$均为渐近稳定。很容易表明，如果满足
\begin{equation}\label{eq:7.4.49}
\min\left\{\lambda_{\min}\{K_{Lp}\},\lambda_{\min}\{K_{Lv}\},\lambda_{\min}\{K_{Fv}\},\lambda_{\min}\{K_{Fp}\}\right\} > \frac{2}{m_1} + 1
\end{equation}
给出的充分条件，式\eqref{eq:7.4.42}和式\eqref{eq:7.4.48}给出的条件总是满足的。

应该注意的是，式\eqref{eq:7.4.39}给出的校正控制器取决于$u_L$、$\dot{u}_L$、$\ddot{u}_L$、$q_m$及$\dot{q}_m$。同样，似乎这个控制需要测量$\ddot{q}$及其导数；然而，由于我们假设精确知道动力学模型，我们可以使用这些信息来消除测量$\ddot{q}$及其导数的需要。也就是说，$\dot{u}_L$可以写成
\begin{equation}\label{eq:7.4.50}
\dot{u}_L = \dot{M}(q)(\ddot{q}_d + K_{Lv}\dot{e} + K_{Lp}e) + \dot{\mathbf{N}}(q,\dot{q}) + M(q)\left[\frac{\mathrm{d}}{\mathrm{d}t}\ddot{q}_d + K_{Lv}(\ddot{q}_d - \ddot{q}) + K_{Lp}\dot{e}\right]
\end{equation}
其中$\ddot{q}$由式\eqref{eq:7.4.28}得到为
\begin{equation}\label{eq:7.4.51}
\ddot{q} = M^{-1}(q)\bigl[Kq_m - \mathbf{N}(q,\dot{q})\bigr]
\end{equation}
其中$\mathbf{N}(q,\dot{q})$由式\eqref{eq:7.4.30}给出。

将式\eqref{eq:7.4.51}代入式\eqref{eq:7.4.50}，我们看到$\dot{u}_L$可以写成$t$、$q$、$\dot{q}$与$q_m$的函数，因为期望轨迹可以明确地写成时间的函数。我们可以使用以下方程阐述这种函数依赖性
\begin{equation}\label{eq:7.4.52}
\dot{u}_L = f(t,q,\dot{q},q_m)
\end{equation}
其中$f(\cdot)$为式\eqref{eq:7.4.50}右侧给出的$n \times 1$向量函数。我们可以通过对式\eqref{eq:7.4.52}关于时间求导获得$\ddot{u}_L$的函数依赖性
\begin{equation}\label{eq:7.4.53}
\ddot{u}_L = \dot{f}(t,q,\dot{q},q_m) = g(t,q,\dot{q},\ddot{q},q_m,\dot{q}_m)
\end{equation}
其中$g(t,q,\dot{q},\ddot{q},q_m,\dot{q}_m)$为$n \times 1$向量函数。注意式\eqref{eq:7.4.51}可以再次用于消除测量$\ddot{q}$的需要；因此，式\eqref{eq:7.4.50}和式\eqref{eq:7.4.53}中给出的$\dot{u}_L$和$\ddot{u}_L$的表达式仅取决于$q$、$\dot{q}$、$q_m$和$\dot{q}_m$的测量。在控制输入$u_F$处实现的实际控制可以通过进行适当的代入找到。也就是说，式\eqref{eq:7.4.39}给出的校正控制可以写成
\begin{equation}\label{eq:7.4.54}
\begin{aligned}
u_F &= J K^{-1}\ddot{u}_L + B_m(q_m,\dot{q}_m,q) + J K^{-1}K_{Fv}(\dot{u}_L - K\dot{q}_m) \\
&\quad {} + J K^{-1}K_{Fp}\bigl[M(q)(\ddot{q}_d + K_{Lv}\dot{e} + K_{Lp}e) + \mathbf{N}(q,\dot{q}) - K q_m\bigr]
\end{aligned}
\end{equation}
其中$\ddot{u}_L$由式\eqref{eq:7.4.53}给出；方括号内即为$\eta_F=u_L-Kq_m$（见式\eqref{eq:7.4.34}与式\eqref{eq:7.4.36}）。

在仔细检查式\eqref{eq:7.4.50}和\eqref{eq:7.4.53}中给出的$\dot{u}_L$和$\ddot{u}_L$的函数依赖性后，很明显为什么我们假设期望轨迹和连杆动力学足够平滑。具体而言，从式\eqref{eq:7.4.50}和\eqref{eq:7.4.53}我们可以看到，式\eqref{eq:7.4.39}中给出的校正控制器将需要期望轨迹的第一、第二、第三和第四时间导数有界，同时需要连杆动力学的第一和第二导数存在。对期望轨迹和连杆动力学的这些假设确保控制输入保持有界。

上面推导的校正控制器总结在表7.4.2中，并在图7.4.4中描绘。

\textbf{表7.4.2}\quad RLFJ校正控制器

\begin{center}
\small
\begin{tabular}{|l|p{11.2cm}|}
\hline
\textbf{扭矩控制器} & \begin{minipage}[t]{\linewidth}\raggedright\setlength{\lineskip}{6pt}
$\begin{aligned}[t] u_F &= J K^{-1}\ddot{u}_L + B_m(q_m,\dot{q}_m,q) + J K^{-1}K_{Fv}(\dot{u}_L-K\dot{q}_m) \\ &\quad {} + J K^{-1}K_{Fp}\bigl[M(q)(\ddot{q}_d+K_{Lv}\dot{e}+K_{Lp}e)+\mathbf{N}(q,\dot{q})-Kq_m\bigr] \end{aligned}$\par
$\begin{aligned}[t] \dot{u}_L &= \dot{M}(q)(\ddot{q}_d+K_{Lv}\dot{e}+K_{Lp}e)+\dot{\mathbf{N}}(q,\dot{q}) \\ &\quad {} + M(q)\left[\frac{\mathrm{d}}{\mathrm{d}t}\ddot{q}_d+K_{Lv}(\ddot{q}_d-\ddot{q})+K_{Lp}\dot{e}\right] \end{aligned}$\par
$\ddot{q}=M^{-1}(q)\bigl[Kq_m-\mathbf{N}(q,\dot{q})\bigr]$\par
$\ddot{u}_L=\dfrac{\mathrm{d}}{\mathrm{d}t}\dot{u}_L\quad\text{（利用 $\ddot{q}$）}$
\end{minipage} \\
\hline
\textbf{稳定性} & 跟踪误差$e$与$\dot{e}$渐近稳定。 \\
\hline
\textbf{注释} & 控制器需要系统动力学的精确知识；增益矩阵$K_{Lv}$、$K_{Lp}$、$K_{Fv}$与$K_{Fp}$须足够大；期望轨迹须足够光滑。 \\
\hline
\end{tabular}
\end{center}

其中$\dot{u}_L$与$\ddot{u}_L$分别由表中第二式及$\ddot{u}_L=\frac{\mathrm{d}}{\mathrm{d}t}\dot{u}_L$给出，并与式\eqref{eq:7.4.50}、式\eqref{eq:7.4.53}一致；$\ddot{q}$由式\eqref{eq:7.4.28}得到即表中第三式，代入后可消去对$\ddot{q}$的直接测量需求。表中第一式与式\eqref{eq:7.4.39}等价（$\eta_F=u_L-Kq_m$，$\dot{\eta}_F=\dot{u}_L-K\dot{q}_m$）。

我们现在提出一个示例来说明如何使用表7.4.2为RLFJ机器人设计校正控制器。

\textbf{示例7.4--2：单连杆RLFJ机械臂的校正控制器}

我们希望使用表7.4.2为图7.4.5中给出的单连杆柔性关节机器人臂设计和仿真校正控制器。系统的动力学取为
\begin{align}
m L^2 \ddot{q} + m L g \sin q + f_d \dot{q} + K q &= K q_m, \tag{1} \\
J \ddot{q}_m + B \dot{q}_m + K(q_m - q) &= u_F. \tag{2}
\end{align}
（与式\eqref{eq:7.4.28}--\eqref{eq:7.4.30}对应：$M=mL^2$，$N=mLg\sin q+f_d\dot{q}$，$\mathbf{N}=N+Kq$，且$B_m(q_m,\dot{q}_m,q)=B\dot{q}_m+K(q_m-q)$。）
其中$m=1$ kg，$K=10$ N，$B=5$ kg-m/s，$f_d=3$ kg-m/s，$L=1$ m，$g$为重力加速度，$J=0.2$ kg-m$^2$，$q_m$为电机位移（弧度），$u_F$为电机输入扭矩。

假设式(1)和(2)给出的模型精确已知，由式\eqref{eq:7.4.54}可得该校正控制器为
\begin{equation}\tag{3}
\begin{aligned}
u_F &= J K^{-1}\ddot{u}_L + B\dot{q}_m + K(q_m-q) + J K^{-1}K_{Fv}(\dot{u}_L-K\dot{q}_m) \\
&\quad {} + J K^{-1}K_{Fp}\bigl[mL^2(\ddot{q}_d+K_{Lv}\dot{e}+K_{Lp}e)+mLg\sin q+f_d\dot{q}+Kq-Kq_m\bigr],
\end{aligned}
\end{equation}
\begin{equation}\tag{4}
\dot{u}_L = mLg\dot{q}\cos q+f_d\ddot{q}+K\dot{q}+mL^2\left[\frac{\mathrm{d}}{\mathrm{d}t}\ddot{q}_d+K_{Lv}(\ddot{q}_d-\ddot{q})+K_{Lp}\dot{e}\right],
\end{equation}
其中$\ddot{q}$由式(1)得到为
\begin{equation}\tag{5}
\ddot{q}=(mL^2)^{-1}\bigl[Kq_m-mLg\sin q-f_d\dot{q}-Kq\bigr].
\end{equation}
将式(5)代入式(4)得到
\begin{equation}\tag{6}
\begin{aligned}
\dot{u}_L &= mLg\dot{q}\cos q+f_d(mL^2)^{-1}\bigl[Kq_m-mLg\sin q-f_d\dot{q}-Kq\bigr] \\
&\quad {} +K\dot{q}+mL^2\left[\frac{\mathrm{d}}{\mathrm{d}t}\ddot{q}_d+K_{Lv}\left(\ddot{q}_d-(mL^2)^{-1}\bigl[Kq_m-mLg\sin q-f_d\dot{q}-Kq\bigr]\right)+K_{Lp}\dot{e}\right].
\end{aligned}
\end{equation}
再对式(6)关于时间求导以得到$\ddot{u}_L$，可得
\begin{equation}\tag{7}
\begin{aligned}
\ddot{u}_L &= mLg\ddot{q}\cos q-mLg\dot{q}^2\sin q \\
&\quad {} +f_d(mL^2)^{-1}\bigl[K\dot{q}_m-mLg\dot{q}\cos q-f_d\ddot{q}-K\dot{q}\bigr]+K\ddot{q} \\
&\quad {} +mL^2\left[\frac{\mathrm{d}^2}{\mathrm{d}t^2}\ddot{q}_d+K_{Lv}\left(\frac{\mathrm{d}}{\mathrm{d}t}\ddot{q}_d-(mL^2)^{-1}\bigl[K\dot{q}_m-mLg\dot{q}\cos q-f_d\ddot{q}-K\dot{q}\bigr]\right)+K_{Lp}(\ddot{q}_d-\ddot{q})\right].
\end{aligned}
\end{equation}

校正控制器在以下条件下进行了仿真：控制参数、初始条件和期望轨迹由
\begin{equation*}
K_{Lv}=K_{Lp}=K_{Fv}=K_{Fp}=5
\end{equation*}
和
\begin{equation*}
q_d=\sin t,\quad q(0)=\dot{q}(0)=q_m(0)=\dot{q}_m(0)=0
\end{equation*}
给出。

跟踪误差和控制扭矩在图7.4.6中描绘。如图所示，跟踪误差是渐近稳定的。

\section{总结}

在本章中，给出了机器人机械臂控制的几种更先进控制技术的描述。本章的意图是研究减少在线计算的控制器和补偿执行器动力学的控制器。一些当前的研究问题涉及力控制器与高级运动控制器的集成以及相应的数字实现。

\vspace{1cm}
\noindent\textbf{参考文献}

\begin{enumerate}
\item[{[1]}] Corless, M., and G. Leitmann, ``Adaptive control of systems containing uncertain functions and unknown functions with uncertain bounds,'' \textit{J. Optim. Theory Appl.}, Jan. 1983.

\item[{[2]}] Dawson, D.M., Z. Qu, F.L. Lewis, and J.F. Dorsey, ``Robust control for the tracking of robot motion,'' \textit{Int. J. Control}, vol. 52, pp. 581--595, 1990.

\item[{[3]}] Eppinger, S., and W. Seering, ``Introduction to Dynamic models for robot force control,'' \textit{IEEE Control Syst. Mag.}, vol. 7, no. 2, pp. 48--52, Apr. 1987.

\item[{[4]}] Ghorbel, F., and M. Spong, ``Stability analysis of adaptively controlled flexible joint robots,'' \textit{Proc. IEEE Conf. Decision Control}, Honolulu, pp 2538--2544, 1990.

\item[{[5]}] Kokotovic, P.V., H. Khalil, and J. O'Reilly, \textit{Singular Perturbation Methods in Control Analysis and Design}. New York: Academic Press, 1986.

\item[{[6]}] Sadegh, N., and R. Horowitz, ``Stability and robustness analysis of a class of adaptive controllers for robotic manipulators,'' \textit{Int. J. Robot. Res.}, vol. 9, no. 3, pp. 74--92, June 1990.

\item[{[7]}] Sadegh, N., R. Horowitz, W. Kao, and M. Tomizuka, ``A unified approach to the design of adaptive and repetitive controllers for robotic manipulators,'' \textit{Trans. ASME}, vol. 112, pp. 618--629, Dec. 1990.

\item[{[8]}] Spong, M., ``Modeling and control of elastic joint robots,'' \textit{J. Dyn. Syst., Meas. Control}, vol. 109, pp. 310--319, Dec. 1987.

\item[{[9]}] Tarn, T., A. Bejczy, X. Yun, and Z. Li, ``Effect of motor dynamics on nonlinear feedback robot arm control,'' \textit{IEEE Trans. Robot. Autom.}, vol. 7, pp. 114--122, Feb. 1991.

\item[{[10]}] Taylor, D., ``Composite control of direct-drive robots,'' \textit{Proc. IEEE Conf. Decision Control}, pp. 1670--1675, Dec. 1989.
\end{enumerate}

\vspace{1cm}
\noindent\textbf{习题}

\textbf{第7.2节}

\begin{enumerate}
\item[7.2--1] 如果滤波跟踪误差定义为
\[ r = \dot{e} + \Lambda e \]
其中$\Lambda$是正定对角矩阵，说明如何修改DCAL稳定性分析。

\item[7.2--2] 为第3章中给出的二连杆极坐标机器人臂设计和仿真表7.2.1中给出的DCAL控制器。（忽略该机器人有棱柱形连杆的事实。）

\item[7.2--3] DCAL稳定性分析（以及因此控制器本身）能否修改为考虑问题7.2--2中给出的棱柱形连杆机器人？如果可以，说明如何修改。

\item[7.2--4] 如果学习项$\hat{u}_d(t)$被迫保持在先验界限内
\[ \hat{u}_{d\min} \leq \hat{u}_{di}(t) \leq \hat{u}_{d\max} \]
其中下标$i$用于表示$n \times 1$向量$\hat{u}_d(t)$的第$i$个分量，$\hat{u}_{d\min}$、$\hat{u}_{d\max}$是标量常数，说明如何修改RCL稳定性分析以表明速度跟踪误差是渐近稳定的。

\item[7.2--5] 为第3章中给出的二连杆极坐标机器人臂设计和仿真表7.2.2中给出的RCL控制器。（忽略该机器人有棱柱形连杆的事实。）

\item[7.2--6] RCL稳定性分析（以及因此控制器本身）能否修改为考虑问题7.2--5中给出的棱柱形连杆机器人？如果可以，说明如何修改。
\end{enumerate}

\textbf{第7.3节}

\begin{enumerate}
\item[7.3--1] 为第3章中给出的二连杆极坐标机器人臂设计和仿真表7.3.1中给出的自适应鲁棒控制器。（忽略该机器人有棱柱形连杆的事实。）

\item[7.3--2] 自适应鲁棒控制器稳定性分析（以及因此控制器本身）能否修改为考虑问题7.3--1中给出的棱柱形连杆？如果可以，说明如何修改。

\item[7.3--3] 说明如何使用Barbalat引理（见第2章）修改自适应鲁棒控制器的稳定性分析以保证速度跟踪误差也是渐近稳定的。
\end{enumerate}

\textbf{第7.4节}

\begin{enumerate}
\item[7.4--1] 对于常数、对称、正定、$n \times n$矩阵$A$，表明
\begin{enumerate}
\item[(a)] $\|A\|_{i2} = \lambda_{\max}\{A\}$
\item[(b)] $\lambda_{\min}\{A^{-1}\} = \frac{1}{\lambda_{\max}\{A\}}$
\end{enumerate}

\item[7.4--2] 为第3章中给出的二连杆旋转机器人臂设计和仿真表7.4.1中给出的RLED校正控制器。假设两个电机都可以建模为示例7.4--1中给出的电机。

\item[7.4--3] 对于常数、对称、正定、$n \times n$矩阵$A$，表明
\[ A^{-1} = \begin{bmatrix} O_{n \times n} & I_{n \times n} \\ -A_{21} & -A_{22} \end{bmatrix} \]
其中$O_{n \times n}$是$n \times n$零矩阵。

\item[7.4--4] 为第3章中给出的二连杆旋转机器人臂设计和仿真表7.4.2中给出的RLFJ校正控制器。假设两个关节都可以建模为类似于示例7.4--2中给出的关节。
\end{enumerate}

\ifstandalone
  \end{document}
\fi
